% !TEX program = xelatex
% 大物实验 · 实验 3.12/3.13 分光计 · 一本通（五环闭环结构）
% 结构：§P 总图 → §0 基础桥 → J0 调节与读数 → J1 顶角 → J2 最小偏向角与折射率
%       → J3 光栅衍射角 → J4 光栅常量与角色散 → 附录 A/B/C/D
\documentclass[11pt]{ctexart}
%=====================================================================
% theme-violet.tex —— 讲义精讲 · 紫花主题
% 用法：\documentclass[11pt]{ctexart}
%       \input{assets/theme-violet}
%       \begin{document} ... \end{document}
% 只用 ctexart 自带能力 + xcolor/enumitem/fancyhdr/pifont/amsmath，
% 不依赖 tikz/tcolorbox，保证 MiKTeX 可直接编译。
%=====================================================================

\usepackage[a4paper,left=2.0cm,right=2.0cm,top=2.1cm,bottom=2.3cm,%
            headsep=12pt,footskip=22pt]{geometry}
\usepackage{amsmath,amssymb}
\usepackage{needspace}
\usepackage{xcolor}
\usepackage{array}
\usepackage{enumitem}
\usepackage{pifont}
\usepackage{fancyhdr}
\usepackage[hidelinks]{hyperref}

%---------------------------------------------------------------- 配色
\definecolor{violet}{RGB}{107,78,143}      % 主色：紫
\definecolor{violetdeep}{RGB}{63,44,92}    % 深紫：标题/强调
\definecolor{petal}{RGB}{183,159,208}      % 花瓣：装饰线、花饰
\definecolor{bandviolet}{RGB}{247,243,251} % 浅紫底：档案/提示框
\definecolor{bandcream}{RGB}{252,249,242}  % 米白底：易错/注意框
\definecolor{rulepurple}{RGB}{217,207,230} % 分隔线
\definecolor{inkgray}{RGB}{92,96,110}      % 次要文字

%---------------------------------------------------------------- 字体
% 有 Noto CJK 用 Noto，没有就退回 ctex 默认（Windows 下为宋体/黑体）
\IfFontExistsTF{Noto Serif CJK SC}{\setCJKmainfont{Noto Serif CJK SC}}{}
\IfFontExistsTF{Noto Sans CJK SC}{\setCJKsansfont{Noto Sans CJK SC}}{}

%---------------------------------------------------------------- 版式
\setlist{nosep,leftmargin=1.9em}
\linespread{1.28}
\setlength{\parskip}{3.5pt}
\setlength{\parindent}{0pt}
\setlength{\tabcolsep}{5pt}
\renewcommand{\arraystretch}{1.28}

%---------------------------------------------------------------- 花饰
% \ding{95} = 五瓣小花，\ding{96} = 星形花；都来自 pifont，必定存在
\newcommand{\orn}{\textcolor{petal}{\ding{95}}}          % 行内小标
\newcommand{\ornstar}{\textcolor{petal}{\ding{96}}}      % 星形花
\newcommand{\flowerrule}{%                                % 居中花饰分隔线
  \par\vspace{7pt}%
  \noindent\makebox[\linewidth]{\color{petal}%
    \rule{0.20\linewidth}{0.6pt}\hspace{8pt}{\small\ding{95}}\hspace{4pt}%
    \raisebox{1.1pt}{\tiny\ding{95}}\hspace{4pt}{\small\ding{95}}\hspace{8pt}%
    \rule{0.20\linewidth}{0.6pt}}%
  \par\vspace{7pt}}

%---------------------------------------------------------------- 标题与节头
\newcommand{\doctitle}[3]{%   #1 主标题  #2 副标题  #3 元信息（可用 \\ 分行）
  \begin{center}
    {\color{petal}\small \orn\enspace #2\enspace \ornstar}\\[7pt]
    {\LARGE\bfseries\color{violetdeep} #1}\\[10pt]
    {\color{petal}\rule{0.55\linewidth}{1pt}}\\[4pt]
    {\small\color{inkgray} #3}
  \end{center}\vspace{2pt}}

\newcommand{\hdr}[2]{%   #1 章节号（如 §0）  #2 标题
  \par\vspace{12pt}%
  \noindent\colorbox{violetdeep}{\color{white}\bfseries\small\;#1\;}\hspace{7pt}%
  {\color{violetdeep}\large\bfseries #2}\par\nobreak
  \vspace{3pt}{\color{petal}\hrule height 1.2pt}\vspace{6pt}}

\newcommand{\subhdr}[1]{%
  \par\vspace{6pt}%
  \noindent{\color{violet}\bfseries #1}\par\vspace{1pt}}

%---------------------------------------------------------------- 提示框
\newcommand{\callout}[1]{% 浅紫底：讲义档案、读法说明
  \par\medskip\noindent
  \fcolorbox{rulepurple}{bandviolet}{\parbox{\dimexpr\linewidth-2\fboxsep-2\fboxrule}{\vspace{2pt}#1\vspace{2pt}}}%
  \par\medskip}

\newcommand{\warnbox}[1]{% 米白底：易错、边界条件
  \par\medskip\noindent
  \fcolorbox{rulepurple}{bandcream}{\parbox{\dimexpr\linewidth-2\fboxsep-2\fboxrule}{\vspace{2pt}#1\vspace{2pt}}}%
  \par\medskip}

\newcommand{\gapitem}[2]{% #1 缺口类型（紫） #2 内容
  \par\vspace{2pt}\noindent
  \hangindent=1.4em\hangafter=1
  {\color{petal}\small\ding{96}}\enspace{\bfseries\color{violetdeep}#1}\enspace #2\par}

% 题干块（\qstem）：每道例题／作业的「详解」开头必须先复述题干与出处。
% 起因：2026-10-25 第五章讲义把七道作业的题面集中放在一页、详解另起一页，
% 读者翻到「第 4 题」时看不到条件，等于一道没有题干的题目。
% 规则：详解块内自带题干，读者不翻页就能读懂题在问什么。
\newcommand{\qstem}[1]{%
  \par\medskip\noindent
  \fcolorbox{violet}{bandviolet}{%
    \parbox{\dimexpr\linewidth-2\fboxsep-2\fboxrule}{%
      \vspace{2pt}%
      {\small\bfseries\color{violet}题干}\hspace{0.8em}#1%
      \vspace{2pt}}}%
  \par\medskip}

% v2: 保护题干与开头，长题可跨页，续页补上下文。
\newcommand{\qbreak}{\needspace{6\baselineskip}}
% 独立练习：ID、来源、题干、考点、答案标签（通常 ans:ID）。
\newcommand{\practice}[5]{%
  \qbreak\subhdr{练习 #1}\qstem{（#2）#3}%
  \textbf{考点}：#4\par
  详解见附录第~\pageref{#5}~页。\par\medskip}
% 附录开头：ID、来源、条件摘要；自动产生 ans:ID 标签。
\newcommand{\answerstart}[3]{%
  \qbreak\subhdr{#1\quad 详解}\label{ans:#1}\qstem{（#2）#3}}
\newcommand{\answercontinue}[2]{%
  \needspace{5\baselineskip}\subhdr{#1\quad 续解}%
  {\small\color{inkgray}本页接续：#2}\par\medskip}
% 延迟提供默认值，让旧讲义可在导言区保留自己的 \newcommand{\step}。
\AtBeginDocument{\providecommand{\step}[1]{\needspace{3\baselineskip}\par\medskip\textbf{#1}\par}}

% 前提旁注：只放在「需要理解的那一步」后面，给陈述 + 本题代入，不重证。
\newcommand{\aside}[1]{%
  \par\vspace{1pt}\noindent
  \hangindent=1.6em\hangafter=1
  {\small\color{inkgray}%
    \hspace*{0.6em}{\color{petal}$\triangleright$}\hspace{2pt}#1}%
  \par\vspace{2pt}}

%---------------------------------------------------------------- 表格
% 本机 MiKTeX 的 colortbl 与 array 版本不匹配（\rowcolor 会报错），
% 故表格统一用「上粗紫线 + 表头 + 中细线」的三线式，不依赖 colortbl。
\newcommand{\tabtop}{\par\vspace{3pt}{\color{violet}\hrule height 1.2pt}\vspace{1pt}}
\newcommand{\tabmid}{\par\vspace{1pt}{\color{petal}\hrule height 0.5pt}\vspace{1pt}}
\newcommand{\tabbot}{\par\vspace{3pt}{\color{rulepurple}\hrule height 0.5pt}\vspace{8pt}}
\newcommand{\thead}[1]{{\bfseries\color{violetdeep}#1}}

%---------------------------------------------------------------- 页眉页脚
\pagestyle{fancy}
\fancyhf{}
\renewcommand{\headrulewidth}{0.4pt}
\renewcommand{\headrule}{\hbox to\headwidth{\color{rulepurple}\leaders\hrule height \headrulewidth\hfill}}
\fancyhead[L]{\small\color{inkgray}\songti \doctag}
\fancyhead[R]{\small\color{inkgray}\songti \coursetag}
\fancyfoot[C]{\small\color{inkgray}\orn\enspace{\color{violet}\thepage}\enspace\orn}

% 每篇讲义在导言区设这两个值
\newcommand{\doctag}{}
\newcommand{\coursetag}{}
\newcommand{\settags}[2]{\renewcommand{\doctag}{#1}\renewcommand{\coursetag}{#2}}

%---------------------------------------------------------------- 数学记号
\newcommand{\vect}[1]{\boldsymbol{#1}}
\newcommand{\dd}{\,\mathrm{d}}
\newcommand{\ddx}{\frac{\mathrm{d}}{\mathrm{d}x}}
\addtolength{\jot}{2pt}



\settags{分光计 3.12/3.13 一本通}{大物实验}
\usepackage{tikz}
\usepackage{graphicx}
\usepackage{booktabs}
\usepackage{float}
\usepackage{needspace}
\usetikzlibrary{arrows.meta,calc,positioning,decorations.pathmorphing}
\setlength{\headheight}{22pt}

% 圆圈数字归入 CJK 字符类，避免 xeCJK 去 Latin 字体里找
\xeCJKDeclareCharClass{CJK}{"2460->"2473}

\newcommand{\blank}{\textcolor{black!45}{\rule[-2pt]{1.05cm}{0.4pt}}}
\newcommand{\blankcell}{\rule{0pt}{16pt}\textcolor{black!45}{\rule[-3pt]{1.35cm}{0.4pt}}}
\newcommand{\rowstrut}{\rule{0pt}{17pt}}
\newcommand{\tick}{$\square$}

\tikzset{
  comp/.style={fill=white,draw=black,thick},
  arr/.style={->,>=Stealth,thick},
  arrd/.style={->,>=Stealth,thick,dashed},
  lab/.style={font=\small},
  dimlab/.style={font=\footnotesize\color{inkgray}}
}

\begin{document}

\doctitle{分光计：一台仪器，四个数}
{实验 3.12/3.13 一本通 $\cdot$ 调节 $\to$ 顶角 $\to$ 折射率 $\to$ 光栅常量与角色散}
{依据：Exp.12 讲义 15 页 + Preview for Experiment 12 三页 $\cdot$ 2026-11-14（周六）13:30--17:30 $\cdot$ 综合楼 308 $\cdot$ 王淑芬}

%=====================================================================
\hdr{§P}{四小时一张图：五个环节，一条链}

这场实验不是「一个实验」，是\textbf{一台仪器、四个测量}。四个测量共用同一套调节，只是往载物台上换东西；而它们的数还互相喂：顶角 $\alpha$ 是折射率公式的输入，光栅的衍射角 $\varphi_k$ 是光栅常量的输入。所以本书按\textbf{环节}分，不按「原理 / 操作 / 数据」分——每一环从「要交出什么」讲到「怎么写进报告」，中间不再跳页。

\subhdr{P.1 这一场要交出的五个数}

\begin{center}
\begin{minipage}{\linewidth}
\begin{tabular}{@{}p{0.9cm}p{4.6cm}p{3.6cm}p{5.4cm}@{}}
\toprule
\thead{环节} & \thead{交出什么} & \thead{交给谁} & \thead{合格判据} \\
\midrule
J0 & 一台可信的仪器 & 后面全部环节 & 两面都重合、像清晰无视差、双游标恒差 $180^{\circ}$ \\
J1 & 顶角 $\alpha$ & J2 的公式 & $60^{\circ}\pm5'$（等边棱镜） \\
J2 & $\delta_{\min}$ 与折射率 $n$ & 报告 & $n$ 落在 $1.45$--$1.65$ \\
J3 & 一、二级衍射角 $\varphi_1,\varphi_2$ & J4 & $\sin\varphi_2/\sin\varphi_1\approx2$ \\
J4 & 光栅常量 $d$、线数 $N$、角色散 $D$ & 报告 & $N$ 接近整数；$D\approx1'/\mathrm{nm}$ \\
\bottomrule
\end{tabular}
\end{minipage}
\end{center}

\textbf{依赖链}：J0 是地基，后面两条支路并联。

\[
\text{J0}\ \longrightarrow\
\begin{cases}
\text{J1}\ (\alpha)\ \longrightarrow\ \text{J2}\ (\delta_{\min}\to n)\\[2pt]
\text{J3}\ (\varphi_k)\ \longrightarrow\ \text{J4}\ (d,\,N,\,D)
\end{cases}
\]

两条支路\textbf{互不需要对方的数}，所以现场哪条先做都行；但如果仪器没调好，两条会一起偏，而且偏的是系统误差，事后平均不回来。

\subhdr{P.2 每环固定六段}

每一环都是同一套六段，位置固定，翻到哪一环都按同样的顺序读：

\begin{enumerate}[label=\textbf{\arabic*},leftmargin=2.2em]
\item \textbf{这一环交出什么}——产出、判据、用时；
\item \textbf{为什么这样能测}——只讲这一环要用的原理，从你已经有的光学基础接上；
\item \textbf{手怎么动}——编号步骤，每步写清「动哪个件 / 看到什么 / 什么算到位」；
\item \textbf{当场记这些数}——表就放在这里，栏目含义、读几组、什么时候锁紧；
\item \textbf{当场算}——表 $\to$ 公式 $\to$ 代入 $\to$ 自洽判据，做完就知道对不对；
\item \textbf{写进报告的一句话}——方法、数据、结果、误差来源。
\end{enumerate}

\subhdr{P.3 四小时怎么分}

\begin{center}
\begin{minipage}{\linewidth}
\begin{tabular}{@{}p{2.5cm}p{1.3cm}p{9.6cm}@{}}
\toprule
\thead{时段} & \thead{用时} & \thead{做什么（对应本文章节）} \\
\midrule
13:30--13:50 & 20 min & 认零件、目测粗调、望远镜聚焦无穷远（J0.3 步骤 1--2） \\
13:50--14:25 & 35 min & 各半法调到两面重合；准直管两步；读数练 2 次（J0.3 步骤 3--5、J0.5） \\
14:25--14:50 & 25 min & 棱镜摆位、测顶角 3 组（J1） \\
14:50--15:30 & 40 min & 最小偏向角：两个方向各找折返点、各 3 组（J2） \\
15:30--15:45 & 15 min & 换光栅、对零级、调 $b_1b_2$、调等高（J3.3） \\
15:45--16:35 & 50 min & 测紫/绿 $\pm1,\pm2$ 共 8 个方位，每个读两游标（J3.4） \\
16:35--17:05 & 30 min & 现场粗算：$\alpha$、$\delta_{\min}$、$n$、$\varphi_k$、$d$、$D$ 各算一遍（J4.5） \\
17:05--17:30 & 25 min & 补测不自洽的那个数、签字、清点器材 \\
\bottomrule
\end{tabular}
\end{minipage}
\end{center}

\callout{\textbf{只有一句话值得现在记住}：前 50 分钟全花在调仪器上不亏。仪器调到「两面都重合、像清晰无视差、每组都读两个游标」，剩下的事就是把像对准叉丝、把两个数抄下来。这三件事缺任何一件，后面的数据不可救；其余的（少测一组、二级没测）都只是精度损失。}

%=====================================================================
\hdr{§0}{你已经有一半了：三条基础桥}

这门实验最容易劝退人的地方，是它看起来要一整套光学。其实\textbf{你只需要三件已经学过的东西}，加一个新概念。

\subhdr{0.1 这门实验从你的光学课里借什么}

\begin{center}
\begin{minipage}{\linewidth}
\begin{tabular}{@{}p{4.2cm}p{4.8cm}p{6.0cm}@{}}
\toprule
\thead{分光计要用的} & \thead{你已经有的} & \thead{这次只加这一层} \\
\midrule
折射定律、棱镜把光掰弯 & 第一章第一课：反射折射四定律、棱镜、全反射 & 偏向角 $\delta=i+i'-\alpha$，与「什么时候最小」\\
焦面与平行光 & 第一章第三、四课：薄透镜、望远镜 & 分划板放在焦面 $\Rightarrow$ 能看见自己的像\\
光程差与双缝干涉 & 第五章：杨氏双缝、多光束 & $d\sin\varphi=k\lambda$，光栅就是很多条缝\\
折射率随波长变 & 第一章、第四章：界面与折射率 & 角色散 $D=\mathrm{d}\varphi/\mathrm{d}\lambda$ 的定义\\
\bottomrule
\end{tabular}
\end{minipage}
\end{center}

\textbf{你真的要新学的只有一个概念}：一束平行光斜着穿过一排等间距的缝时，相邻两条缝的光「多走了多少路」正好是 $d\sin\varphi$；这个多走的量是波长的整数倍，那个方向上就亮。这一句话就是光栅方程的全部来历，J3 里会用一整节把它讲透。

\callout{\textbf{光学课没讲到的部分，本文按「没学过」处理。}凡是用到衍射强度分布、菲涅耳/夫琅禾费积分的地方，本文一概不用；光栅只从双缝的光程差推，不引入新的波动理论。你在第五章学过的「光程差决定明暗」这一条，足够支撑全书。}

\subhdr{0.2 记号表（全书统一）}

\begin{center}
\begin{minipage}{\linewidth}
\begin{tabular}{@{}p{1.5cm}p{7.6cm}p{2.2cm}@{}}
\toprule
\thead{符号} & \thead{含义} & \thead{单位} \\
\midrule
$\alpha$ & 三棱镜顶角（两折射面的夹角） & $^{\circ}$ \\
$\delta,\ \delta_{\min}$ & 偏向角、最小偏向角 & $^{\circ}$ \\
$i,\ i'$ & 两个面上的入射角 & $^{\circ}$ \\
$r,\ r'$ & 两个面上的折射角 & $^{\circ}$ \\
$n$ & 三棱镜玻璃对该波长的折射率 & --- \\
$f,\ a$ & 望远镜物镜焦距；透光「十」到中心横丝的距离 & mm \\
$\phi_1,\phi_1'$ & 望远镜在位置 I 时\textbf{两个游标}的读数 & $^{\circ}\,'$ \\
$\phi_2,\phi_2'$ & 望远镜在位置 II 时两个游标的读数 & $^{\circ}\,'$ \\
$\theta_{\pm k},\theta_{\pm k}'$ & 第 $\pm k$ 级谱线在两个游标上的读数 & $^{\circ}\,'$ \\
$\varphi_k$ & 第 $k$ 级衍射角（两条光线的夹角） & $^{\circ}$ \\
$d,\ N$ & 光栅常量（相邻缝距）；刻线密度 & nm；线/mm \\
$D$ & 角色散 $\mathrm{d}\varphi/\mathrm{d}\lambda$ & $'/\mathrm{nm}$ \\
\bottomrule
\end{tabular}
\end{minipage}
\end{center}

\warnbox{\textbf{两处一开始就要分清，否则每一步都会代错。}\\
\textbf{(1) $\phi$ 与 $\varphi$ 不是一回事。}$\phi$ 是望远镜在度盘上的\textbf{方位读数}，$\varphi_k$ 是\textbf{两条光线的夹角}（衍射角）。读数必须相减才成为角。\\
\textbf{(2) 带撇与不带撇。}$\phi_1$ 与 $\phi_1'$ 是\textbf{同一次}望远镜位置上的\textbf{两个游标}（相差约 $180^{\circ}$），不是两次测量。算角度时 $|\phi_2-\phi_1|$ 与 $|\phi_2'-\phi_1'|$ 要各自算差、再相加。}

\subhdr{0.3 还没开始测就能知道的四个数}

\begin{center}
\begin{minipage}{\linewidth}
\begin{tabular}{@{}p{6.4cm}p{7.8cm}@{}}
\toprule
\thead{关系} & \thead{现场怎么用} \\
\midrule
等边棱镜 $\alpha=60^{\circ}$，两反射像夹角 $120^{\circ}$ & 望远镜从一个像转到另一个应转过约 $120^{\circ}$；差很多是棱镜没放正 \\
玻璃 $n\approx1.5$（$\alpha=60^{\circ}$）$\Rightarrow\delta_{\min}\approx37^{\circ}$ & 两位置读数差应约 $74^{\circ}$；算出 $n$ 超出 $1.45$--$1.65$ 就回去查读数 \\
$300$ 线/mm、一级绿线 $\Rightarrow\varphi\approx9.4^{\circ}$ & 一级在 $9^{\circ}$ 出头、二级在 $19^{\circ}$；一级只有 $3^{\circ}$ 说明是 $100$ 线/mm \\
双游标读数恒差 $180^{\circ}00'$（$\pm5'$） & 每读一次顺手对一眼，偏了就是读错游标或漏了 $360^{\circ}$ 折回 \\
\bottomrule
\end{tabular}
\end{minipage}
\end{center}

第三条最有用。汞灯绿线 $\lambda=546.1\ \mathrm{nm}$，若光栅为 $300$ 线/mm，则 $d=1/300\ \mathrm{mm}=3.333\ \mu\mathrm{m}$：

\[
\sin\varphi_1=\frac{\lambda}{d}=\frac{546.1\ \mathrm{nm}}{3333\ \mathrm{nm}}=0.16383
\quad\Longrightarrow\quad
\varphi_1=9.43^{\circ},\qquad \varphi_2=19.13^{\circ}.
\]

所以还没开始测，你就知道一级谱线该出现在零级两边各 $9.4^{\circ}$ 处。转了半圈还找不到，问题不在计算，在光栅没调好（J3.3）。

%=====================================================================
\hdr{J0}{把仪器调到能用：五步，与怎么读数}

\subhdr{① 这一环交出什么}

不交数，交一个\textbf{状态}：三条轴摆正、望远镜聚焦无穷远、游标读数可信。这一环的产物是后面四个测量的前提，用时约 55 分钟。

三条判据，缺一条后面全部带系统误差：
\begin{enumerate}[label=\textbf{(\arabic*)}]
\item 平面镜\textbf{正反两个面}反射的绿十字都与上叉丝重合；
\item 狭缝像\textbf{清晰、竖直、居中}，且眼睛移动时无视差；
\item 每次读数都读两个游标，两值恒差 $180^{\circ}$。
\end{enumerate}

\subhdr{② 为什么这样能测：三条轴与一个自准直}

讲义 p8 把「调好」写成三条要求。翻成一句几何话，后面五步就都有归宿：

\callout{\textbf{望远镜光轴、准直管光轴、仪器主轴三者两两垂直；待测元件的光学面平行于主轴。}}

为什么非要这样：所有角度公式（$\alpha$、$\delta_{\min}$、$d\sin\varphi=k\lambda$）都是在\textbf{垂直于主轴的那个平面}里推的。光一旦离开这个平面，你读到的角就不是公式里的角。这个平面叫\textbf{主截面}（main section），口试问的就是它。

\textbf{自准直为什么能看见自己的像。}望远镜不是普通「物镜 + 目镜」：物镜与目镜之间插了一块\textbf{分划板}，分划板前有一块全反射小棱镜，把绿色 LED 的光引到分划板上那个透光的「十」上——于是「十」变成一枚\textbf{发光的物}，位置就在分划板上。接下来取决于分划板离物镜多远，分三步说：

\begin{enumerate}[label=\textbf{(\arabic*)}]
\item 分划板正好在物镜焦面上时，从「十」发出的光经物镜后成为\textbf{平行光}；但「十」不在光轴上（在轴下方 $a$），所以这束平行光与光轴成角 $\theta\approx a/f$，方向朝上。
\item 平面镜的法线沿光轴，反射只把\textbf{沿轴的分量}反过来，垂直分量不变；返回的仍是平行光，但倾斜到\textbf{另一侧}。
\item 这束平行光再穿过物镜，被聚焦回焦面，落在\textbf{中心横丝上方} $f\tan\theta=a$ 处——正好是\textbf{上叉丝}。
\end{enumerate}

一句话：\textbf{像与物关于光轴对称}。物在轴下 $a$，像就在轴上 $a$。所以判据是「绿十字与\textbf{上叉丝}重合」，不是「像回到『十』身上」。

\begin{figure}[H]
\centering
\begin{tikzpicture}[>=Stealth,line join=round,scale=0.92,transform shape]
  \draw[dash dot,color=black!45] (-1.0,0) -- (8.8,0);
  % 分划板
  \draw[thick] (0,-1.9) -- (0,1.9);
  \draw[thick] (-0.24,0.5) -- (0.24,0.5);
  \draw[thick] (-0.36,0) -- (0.36,0);
  \draw[thick,color=violet] (-0.13,-0.5) -- (0.13,-0.5);
  \draw[thick,color=violet] (0,-0.63) -- (0,-0.37);
  \node[font=\footnotesize,anchor=east,color=violetdeep] at (-0.17,0.50){上叉丝};
  \node[font=\footnotesize,anchor=east] at (-0.17,0){中心横丝};
  \node[font=\footnotesize,anchor=east,color=violet] at (-0.17,-0.50){透光「十」};
  \node[font=\footnotesize] at (0,-2.10){分划板（焦面）};
  % 物镜
  \draw[thick] (3,-1.9) to[bend left=16] (3,1.9);
  \draw[thick] (3,-1.9) to[bend right=16] (3,1.9);
  \node[font=\footnotesize] at (3,-2.10){物镜（焦距 $f$）};
  % 平面镜
  \draw[thick] (8,-1.8) -- (8,2.2);
  \foreach \y in {-1.7,-1.4,...,2.1}{\draw[black!50] (8,\y) -- (8.24,\y-0.24);}
  \node[font=\footnotesize,anchor=west] at (8.10,2.34){平面镜};
  % 光线
  \draw[arr,color=violetdeep] (0,-0.5) -- (3,0);
  \draw[arr,color=violetdeep] (3,0) -- (7.96,0.827);
  \draw[arrd,color=black!70] (7.96,0.827) -- (3,1.6667);
  \draw[arrd,color=black!70] (3,1.6667) -- (0.04,0.507);
  \fill[violetdeep] (0,0.5) circle (0.055);
  \fill[black!70] (8,0.8333) circle (0.045);
  \draw[black!55] (4.70,0) arc[start angle=0,end angle=9.46,radius=0.45];
  \node[font=\footnotesize] at (4.96,0.19){$\theta$};
  \node[font=\footnotesize,anchor=north] at (5.10,-0.42){出射：平行光};
  \node[font=\footnotesize,anchor=south] at (5.60,1.78){返回：平行光};
\end{tikzpicture}
\caption{为什么像落在上叉丝（示意，角度已放大；实际 $\theta=a/f\approx0.6^{\circ}$）。紫实线＝出射，灰虚线＝返回。物点「十」在轴下 $a$，出射平行光上倾；反射后向另一侧倾斜，经物镜落在轴上 $+a$，即上叉丝。}
\end{figure}

\textbf{为什么用自准直而不用眼睛看「对准了」}：镜面偏 $\gamma$，返回光偏 $2\gamma$，像点移开 $2f\gamma$。取 $f=200\ \mathrm{mm}$、$\gamma=1'$：像移约 $0.12\ \mathrm{mm}$，肉眼看得见。角度被\textbf{放大一倍}，这就是它比「大致对准」准得多的原因。

\subhdr{③ 手怎么动：五步}

\textbf{顺序不能换}：步骤 3 用平面镜，步骤 4 用狭缝；后者要望远镜先能看平行光。

\subhdr{步骤 1：目测粗调（5 分钟，决定后面顺不顺）}

\textbf{目标}：两筒大致等高对望，载物台大致水平。\quad
\textbf{动作}：调 12（望远镜光轴竖直）与 27（准直管光轴竖直），眼睛直接沿筒身瞄，让两筒轴线大致在一条直线上；再调载物台三个调平螺钉，使台面大致水平。\quad
\textbf{判据}：把平面镜放上台，镜面大致竖直、反射面朝望远镜。\\
\textbf{为什么不能省}：细调螺钉的行程只有几度，而绿十字像要落进视场（约 $1^{\circ}$ 量级）才看得见。粗调差 $5^{\circ}$，后面怎么调都是白调。

\subhdr{步骤 2：望远镜聚焦无穷远（自准直法）}

\textbf{目标}：分划板落在物镜焦面上。\\
\textbf{动作}：(I) 转 11 号手轮，先把\textbf{叉丝本身}看清；(II) 把双面平面镜放在载物台上，慢慢转载物台让一个反射面朝向望远镜，在视场里找那个\textbf{模糊的绿十字像}（找不到就回步骤 1）；(III) 松 9 号螺钉，前后移目镜筒，让绿像变\textbf{清晰}；(IV) 消视差：眼睛上下左右移，看绿像与叉丝是否相对移动，有就继续微移，直到\textbf{两者一起动}。\\
\textbf{判据}：绿像清晰\textbf{且}无视差。\\
\textbf{为什么}：无视差 $\Leftrightarrow$ 像与叉丝同面 $\Leftrightarrow$ 分划板在焦面 $\Leftrightarrow$ 出射与返回都是平行光。这三句是同一件事的三种说法。

\begin{figure}[H]
\centering
\includegraphics[width=0.46\linewidth]{fig-3.12/fig12-10-mirror.png}
\caption{平面镜在载物台上的两种摆法（讲义 p9 图 12-10）：(a) 镜面平行于一对螺钉连线；(b) 镜面垂直于一对螺钉连线。两种都行，关键是反射面朝望远镜、镜面竖直。$a,b,c$ 是三个调平螺钉。}
\end{figure}

\subhdr{步骤 3：望远镜轴垂直主轴（各半法）}

\textbf{目标}：望远镜光轴 $\perp$ 主轴。\\
\textbf{判据}（讲义 p9）：镜子的\textbf{正反两个面}反射的绿十字都要与上叉丝重合。\\
\textbf{动作}：(1) 当前位置把绿十字调到与上叉丝重合；(2) 载物台转 $180^{\circ}$，找另一个面的像（找不到就用眼睛直接看镜面里像的位置，先把调平螺钉 $a$ 转到能看见）；(3) 看到偏差 $d$：先调\textbf{调平螺钉 $a$} 把距离\textbf{减半}，再调 \textbf{12 号} 让像\textbf{落到}上叉丝上；(4) 再转 $180^{\circ}$ 重复 (3)，直到两面都重合。

\warnbox{\textbf{为什么必须「各半」，不能一次调到位？}两个偏差来源对两个面的影响方式不一样：望远镜轴歪了，两个面的像\textbf{同向}偏（一起高或一起低）；载物台歪了（镜面不垂直主轴），转 $180^{\circ}$ 后误差\textbf{反号}。设同向分量 $v$、反向分量 $u$，则正面偏差 $=v+u$、反面偏差 $=v-u$。调 12 号只改 $v$，调 $a$ 只改 $u$。若在一个面上用 $a$ 一次把偏差消光，等于把 $u$ 全转嫁给另一面而 $v$ 没动；「先减半、再归零」才是每轮把两者同时压一半，两三轮回合就都到位——原文叫它 \emph{asymptotic method}（渐近法）。}

\subhdr{步骤 4：让准直管发出平行光}

\textbf{目标}：狭缝落在准直管透镜的焦面上。\\
\textbf{动作}：取下平面镜，用汞灯照亮狭缝；望远镜转过来对准准直管，从目镜看狭缝像；松 2 号螺钉，前后移狭缝套筒，使像\textbf{清晰、边缘锐、无视差}；再把缝宽调到约 $0.5$~mm。\\
\textbf{判据}：像清晰且无视差。有视差＝狭缝不在焦面＝出来的不是平行光。\\
\textbf{缝宽为什么定 $0.5$~mm}：太宽（$>1$~mm）亮度够，但像成宽带，对准中心全凭主观，两条近谱线还会并成一条；太窄（$<0.2$~mm）谱线细锐却太暗，缝的衍射反而让边缘发虚。$0.5$~mm 是亮度与对准精度的折中。

\subhdr{步骤 5：准直管轴垂直主轴，并让狭缝竖直}

\textbf{动作}：(1) 转狭缝套筒，使像\textbf{平行于横向叉丝}；(2) 调 27 号，把像\textbf{升/降到视场中央}与横向叉丝重合——此时准直管光轴 $\perp$ 主轴；(3) 把套筒\textbf{转 $90^{\circ}$}，让像平行于\textbf{竖直}叉丝；(4) 拧紧 2 号锁住。整个过程不要把套筒前后推动，那会破坏步骤 4 的聚焦。\\
\textbf{判据}：狭缝像竖直、细锐、居中。

\subhdr{步骤 6：这一次一定要记住的螺钉}

\warnbox{\textbf{16 号螺钉决定「谁跟谁一起转」。}松开它，游标盘座与度盘一起转；拧紧它，两者可相对转。测量要的是「望远镜相对度盘转过的角」，所以两个读数必须相对于\textbf{同一个度盘}。规矩：\textbf{一组测量开始后，16 号的状态不许再改}。想锁死望远镜用 17 号，不要动 16 号。}

\subhdr{④ 状态确认清单（调完逐项打勾，不合格就回对应步骤）}

\begin{center}
\begin{minipage}{\linewidth}
\begin{tabular}{@{}p{0.7cm}p{8.6cm}p{4.4cm}@{}}
\toprule
\thead{\tick} & \thead{确认什么} & \thead{不合格回哪} \\
\midrule
\rowstrut \tick & 叉丝本身清楚、边缘锐 & 步骤 2 第 (I) 小步 \\
\rowstrut \tick & 绿十字像清晰、无视差 & 步骤 2 第 (III)(IV) 小步 \\
\rowstrut \tick & 镜面正反两面都与上叉丝重合 & 步骤 3 \\
\rowstrut \tick & 狭缝像清晰、竖直、居中 & 步骤 4、5 \\
\rowstrut \tick & 缝宽约 $0.5$~mm & 步骤 4 \\
\rowstrut \tick & 16 号螺钉状态记下了（松/紧） & 步骤 6 \\
\bottomrule
\end{tabular}
\end{minipage}
\end{center}

\textbf{卡住了查这张表}（讲义没有，按现象倒推）：

\begin{center}
\begin{minipage}{\linewidth}
\begin{tabular}{@{}p{5.0cm}p{4.2cm}p{4.6cm}@{}}
\toprule
\thead{现象} & \thead{最可能的原因} & \thead{回去做哪一步} \\
\midrule
找不到绿色十字像 & 粗调没做好 & 步骤 1 \\
两个面的像都偏高/偏低 & 望远镜轴歪（$v\neq0$） & 步骤 3，调 12 号 \\
一面偏高、一面偏低 & 台面歪（$u\neq0$） & 步骤 3，调 $a$ \\
狭缝像模糊、有视差 & 狭缝不在焦面 & 步骤 4 \\
狭缝像不在视场中央 & 准直管轴未垂直主轴 & 步骤 5，调 27 号 \\
狭缝像斜着 & 狭缝没转正 & 步骤 5，转套筒 \\
读数总差一个固定角 & 望远镜轴不垂直主轴 & 整段重做步骤 3 \\
\bottomrule
\end{tabular}
\end{minipage}
\end{center}

\subhdr{⑤ 怎么读数：三步，与两个自检}

\textbf{分度关系必须记准}：度盘大格 $1^{\circ}$，大格之间有一条半度刻线，所以\textbf{主尺最小分格 $30'$}；游标共 30 格、跨主尺 29 格，每格 $29'$，故\textbf{分度值 $1'$}。

\textbf{一次读数做三步}（示例取自说明书算例）：
\begin{enumerate}[label=\textbf{第 \arabic* 步}]
\item \textbf{找游标零线}：游标上标 \textbf{0} 的那条线，看它落在主尺哪两条刻线之间。
\item \textbf{读主尺}：取零线\textbf{已经越过}的那条主尺刻线（读数更小的那条）。算例中是 $334^{\circ}30'$。
\item \textbf{读游标}：沿游标从 0 往 30 找那条与主尺刻线\textbf{连成一条直线}的线，它是第几格就加几分钟（算例第 $17$ 格，即 $17'$）。
\item \textbf{相加}：$334^{\circ}30'+17'=334^{\circ}47'$。
\end{enumerate}

\textbf{两个自检}：① 结果必须落在第 2 步那两条刻线之间（$334^{\circ}47'$ 在 $334^{\circ}30'$ 与 $335^{\circ}$ 之间，$\checkmark$）；② 游标读数永远小于 $30'$，只改「分」不改「度」。

\warnbox{\textbf{说明书有一处措辞要注意。}原文写「取零线\emph{右边}那条刻线」，按字面会得 $335^{\circ}$，与它自己的算例矛盾。\textbf{以算例为准}：取零线已越过的那条。现场先读一次，用自检 ① 确认方向。}

\textbf{为什么必须读两个游标：偏心差。}度盘中心与主轴中心在加工上不可能重合，于是望远镜真正转过的角 $\phi$ 与盘上读出的角不相等，误差随方位变化：

\[
\varepsilon(\phi)\approx\frac{e}{R}\sin(\phi-\psi),
\]

$e$ 是偏心量、$R$ 是度盘半径、$\psi$ 是偏心方位。两个游标相差 $180^{\circ}$，所以 $\varepsilon(\phi+180^{\circ})=-\varepsilon(\phi)$——\textbf{等大反号}，取平均就消掉：

\[
\theta_1=\phi+\varepsilon,\quad \theta_2=\phi-\varepsilon
\quad\Longrightarrow\quad
\boxed{\;\phi=\frac{\theta_1+\theta_2}{2}\;}
\]

量级：$e=0.02\ \mathrm{mm}$、$R=60\ \mathrm{mm}$ 时单个游标最大误差约 $1.15'\approx1'$，比读数精度还大；平均后只剩二阶项，可忽略。所以\textbf{不读两个游标，等于白测}。

\textbf{三条硬规矩}：① 每读一次方位，两个游标都读；② 先各自算差再平均，不要「先平均读数再相减」（数学等价，现场容易对错游标）；③ 一组测量开始后不许动 16 号。

\textbf{折回规则}（J1、J2 会用）：度盘是 $0$--$360$ 的圆，望远镜实际转的是\textbf{较小的那个角}。若 $|\phi_2-\phi_1|>180^{\circ}$，改用 $360^{\circ}-|\phi_2-\phi_1|$。\textbf{必须对两个游标分别判断}，因为它们可能一个越过 $0^{\circ}$、一个没有。

\subhdr{⑥ 写进报告的一句话}

\callout{
「经自准直法调节，望远镜聚焦于无穷远（绿十字像清晰无视差），平面镜正反两面反射像均与上叉丝重合，准直管出射平行光且狭缝竖直居中，满足三条调节要求。」
}
\warnbox{
\textbf{台前只看这一段}：调好 = ①叉丝清 ②绿像清且无视差 ③镜面两面都重合 ④缝像竖直居中 ⑤缝宽 $0.5$~mm ⑥不动 16 号。\\
读数 = 主尺（越过的那条）$+$ 游标（对齐的第几格）；两个游标都读，差值应 $180^{\circ}\pm5'$。
}
%=====================================================================
\clearpage
\hdr{J1}{测三棱镜顶角：反射法}

\subhdr{① 这一环交出什么}

\begin{center}
\begin{minipage}{\linewidth}
\begin{tabular}{@{}p{3.0cm}p{4.4cm}p{6.6cm}@{}}
\toprule
\thead{产出} & \thead{合格判据} & \thead{交给谁} \\
\midrule
顶角 $\alpha$ 及其不确定度 & 等边棱镜应 $59^{\circ}55'$--$60^{\circ}05'$ & J2 的折射率公式要用 $\alpha$ \\
\bottomrule
\end{tabular}
\end{minipage}
\end{center}

用时约 25 分钟：摆棱镜与调平 10 分钟，三个位置读数 15 分钟。

\subhdr{② 为什么这样能测：两束反射光的夹角是 $2\alpha$}

把一束方向固定的平行光同时照到三棱镜的两个折射面上（让棱镜的顶点对着准直管），两个面各反射一束光。\textbf{这两束反射光的夹角等于 $2\alpha$}。

\textbf{证明}：单个平面镜转过 $\gamma$，反射光线转过 $2\gamma$——这是反射定律的直接推论：入射方向不变时 $k'=k-2(k\cdot n)n$，法线转 $\gamma$，$k'$ 就转 $2\gamma$。现在两个面不是同一面镜子，但它们对同一束入射光的取向差恰好是顶角 $\alpha$，所以两束反射光的取向差是 $2\alpha$。

\warnbox{\textbf{这里有个很容易答错的地方：为什么不是 $2(180^{\circ}-\alpha)$？}棱镜两面的\textbf{外法线}夹角确实是 $180^{\circ}-\alpha$（$\alpha=60^{\circ}$ 时法线夹角 $120^{\circ}$）。但两束反射光的方向差是 $2\times120^{\circ}=240^{\circ}$，而\textbf{两条直线的夹角取小于 $180^{\circ}$ 的那个}，$360^{\circ}-240^{\circ}=120^{\circ}=2\alpha$。所以无论从哪条路算，望远镜要转的都是 $2\alpha=120^{\circ}$。}

\textbf{公式里的 $1/4$ 是两个 $\tfrac12$ 叠起来的}，不必死记：

\begin{enumerate}[label=\textbf{第 \arabic* 个 $\tfrac12$}]
\item \textbf{几何的 $\tfrac12$}：望远镜在位置 I、II 之间转过的角等于 $2\alpha$，所以 $\alpha=\tfrac12\times$ 转角；
\item \textbf{读数的 $\tfrac12$}：转角本身要用两个游标平均，$\text{转角}=\tfrac12(|{\phi_2-\phi_1}|+|{\phi_2'-\phi_1'}|)$。
\end{enumerate}

两式相乘即 $\tfrac14$。看到 $\tfrac14$ 就想「双面镜 $\times$ 双游标」。

\subhdr{③ 手怎么动}

\textbf{第 1 步：摆棱镜。}按讲义 p11 图 12-11 的方位放：让 $AB$、$AC$、$BC$ 分别垂直于三个调平螺钉两两的连线 $a_1a_2$、$a_2a_3$、$a_3a_1$；顶点稍微靠近载物台中心，这样两个面都能被准直管的光照到。

\begin{figure}[H]
\centering
\includegraphics[width=0.40\linewidth]{fig-3.12/fig12-11-prism-stage.png}
\caption{三棱镜在载物台上的方位（讲义 p11 图 12-11）：$AB\perp a_1a_2$、$AC\perp a_2a_3$、$BC\perp a_3a_1$。这样 $a_1$ 只改变 $AB$ 面法线、$a_3$ 只改变 $AC$ 面法线，两个面可以各自独立调平。}
\end{figure}

\textbf{第 2 步：调平棱镜。}用自准直法：先找到 $AB$ 面的绿十字像，调 $a_1$ 让它与上叉丝重合；再转载物台找 $AC$ 面的像，调 $a_3$ 让它重合；反复直到两面都重合。此时\textbf{棱镜主截面垂直主轴}。

\warnbox{\textbf{为什么偏偏不能调 $a_2$}：在这个摆法下 $a_1$ 只改 $AB$ 面法线、$a_3$ 只改 $AC$ 面法线，而 $a_2$ 会\textbf{同时}改动两个面——刚调好的一个面立刻被破坏，两个独立的问题被耦合成一个解不开的。讲义 p11 只用一句「Do not adjust the screw a2.」带过，原因在这里。}

\textbf{第 3 步：位置 I 读数。}把望远镜转到能看见 $AB$ 面反射的狭缝像，用\textbf{竖直叉丝}对准像的中心，锁 17 号，读\textbf{两个游标}得 $\phi_1$ 与 $\phi_1'$。

\textbf{第 4 步：位置 II 读数。}把望远镜转到 $AC$ 面的反射像，同样对准，读 $\phi_2$ 与 $\phi_2'$。一组就是这四个数。

\textbf{第 5 步：重复三组。}每组之间可以重新对一次像，但\textbf{不要动 16 号}。

\textbf{找像的先后}：如果某一面的像死活找不到，先看棱镜是否放得太偏（顶点没朝准直管）、台面高矮是否挡住光路；再看该面的调平螺钉是否把它调出了视场。

\subhdr{④ 当场记这些数}

\begin{center}
\begin{minipage}{\linewidth}
\small
\begin{tabular}{@{}p{1.4cm}p{2.7cm}p{2.7cm}p{2.7cm}p{2.7cm}p{2.0cm}@{}}
\toprule
\thead{次数} & \thead{$\phi_1$（$AB$ 面）} & \thead{$\phi_1'$} & \thead{$\phi_2$（$AC$ 面）} & \thead{$\phi_2'$} & \thead{备注} \\
\midrule
\rowstrut 1 & \blankcell & \blankcell & \blankcell & \blankcell & \\
\rowstrut 2 & \blankcell & \blankcell & \blankcell & \blankcell & \\
\rowstrut 3 & \blankcell & \blankcell & \blankcell & \blankcell & \\
\bottomrule
\end{tabular}
\end{minipage}
\end{center}

\textbf{栏目含义}：$\phi_1$ 与 $\phi_1'$ 是\textbf{同一次}望远镜位置的两个游标（差约 $180^{\circ}$），不是两次测量。若中间动过 16 号，四个数就不再对应同一个度盘，\textbf{整组作废}。

\subhdr{⑤ 当场算}

\begin{enumerate}[label=\textbf{(\arabic*)}]
\item 每行各算两个差：$|\phi_2-\phi_1|$ 与 $|\phi_2'-\phi_1'|$；\textbf{任一个超过 $180^{\circ}$，就用 $360^{\circ}$ 减它}（这条必须在每一行上分别做，不能等平均之后再判断一次）。
\item 三行的两个差各自取平均，得 $\overline{|\Delta_1|}$ 与 $\overline{|\Delta_2|}$。
\item 代公式：
\[
\alpha=\frac14\Bigl(\overline{|\phi_2-\phi_1|}+\overline{|\phi_2'-\phi_1'|}\Bigr).
\]
\end{enumerate}

\textbf{为什么是 $\tfrac14$ 而不是 $\tfrac12$}：望远镜从位置 I 转到位置 II 的角度等于 $2\alpha$，这是唯一的几何事实。两个游标各给这个转角一个估计 $|\Delta_1|$ 与 $|\Delta_2|$，取平均得
\[
\text{转角}=\frac{|\Delta_1|+|\Delta_2|}{2}=2\alpha
\quad\Longrightarrow\quad
\alpha=\frac{|\Delta_1|+|\Delta_2|}{4},
\]
即手册 (12-4)。一个 $\tfrac12$ 来自「转角是顶角的两倍」，另一个 $\tfrac12$ 来自「两个游标取平均」。

（手册把「平均」写在对四个读数上；没有折回时两者结果相同。有折回时以「每行先折回」为准，因为折回是个非线性操作。）

\textbf{例算（示例数据，非实测）}：$\phi_1=100^{\circ}12'$、$\phi_1'=280^{\circ}13'$、$\phi_2=220^{\circ}13'$、$\phi_2'=40^{\circ}14'$。

\[
|\phi_2-\phi_1|=119^{\circ}61'=120^{\circ}01',\qquad
|\phi_2'-\phi_1'|=239^{\circ}59'\xrightarrow{\ \times\ }120^{\circ}01',
\]
\[
\alpha=\frac14\bigl(120^{\circ}01'+120^{\circ}01'\bigr)=60^{\circ}00.5'.
\]

\textbf{自查三条}：① $\alpha$ 应在 $60^{\circ}$ 附近 $5'$ 以内；② 两个差的平均应约 $120^{\circ}$（$=2\alpha$）；③ 两组游标各自算出的 $\alpha$ 相差应在 $2'$ 内，差多了说明某个游标读错。

\begin{figure}[H]
\centering
\includegraphics[width=0.36\linewidth]{fig-3.12/fig12-12-vertex.png}
\caption{反射法测顶角的两个望远镜位置（讲义 p11 图 12-12）：顶点 $A$ 对住准直管，$AB$、$AC$ 两面各反射一束，望远镜在 I、II 两处读 $\phi_1,\phi_1'$ 与 $\phi_2,\phi_2'$。}
\end{figure}

\subhdr{⑥ 写进报告的一句话}

\callout{
「用反射法测得三棱镜顶角 $\alpha=60^{\circ}00.5'$（三组平均）；两反射光束夹角为 $2\alpha=120^{\circ}01'$，与等边棱镜标称值 $60^{\circ}$ 相符，偏差主要来自度盘估读与棱镜主截面未严格垂直主轴。」
}
\warnbox{
\textbf{台前只看这一段}：棱镜按图摆（$AB\perp a_1a_2$、$AC\perp a_2a_3$），只调 $a_1$、$a_3$，\textbf{不碰 $a_2$}；两面都重合后，位置 I 读两游标、位置 II 读两游标；$\alpha=\tfrac14\bigl(|\phi_2-\phi_1|+|\phi_2'-\phi_1'|\bigr)$，任一差超 $180^{\circ}$ 先折回。
}
%=====================================================================
\clearpage
\hdr{J2}{测最小偏向角，进而得折射率}

\subhdr{① 这一环交出什么}

\begin{center}
\begin{minipage}{\linewidth}
\begin{tabular}{@{}p{3.4cm}p{4.0cm}p{6.6cm}@{}}
\toprule
\thead{产出} & \thead{合格判据} & \thead{用谁的数据} \\
\midrule
$\delta_{\min}$ 与 $n$ & $n$ 落在 $1.45$--$1.65$；两位置读数差约 $74^{\circ}$ & $\delta_{\min}$ 本环测；$\alpha$ 用 J1 的 \\
\bottomrule
\end{tabular}
\end{minipage}
\end{center}

用时约 40 分钟：一个方向找折返点 12 分钟，换面再找 12 分钟，读数与复核 16 分钟。

\subhdr{② 为什么这样能测：偏向角、它与顶角的关系、以及「什么时候最小」}

一束单色平行光射到棱镜 $AB$ 面，折射后从 $AC$ 面射出。入射方向与出射方向的夹角叫\textbf{偏向角} $\delta$。三条关系，把每个三角形里的角标清就能推：

\begin{enumerate}[label=\textbf{(\arabic*)}]
\item 两个面的折射定律：$\sin i=n\sin r$，$\quad n\sin r'=\sin i'$。
\item 顶角关系：两个法线夹角是 $180^{\circ}-\alpha$，三角形 $AOO'$ 的内角和给出 $r+r'=\alpha$。
\item 偏向角：光线在 $AB$ 面偏了 $i-r$，在 $AC$ 面又偏了 $i'-r'$，两次同向叠加：$\delta=(i-r)+(i'-r')=i+i'-\alpha$。
\end{enumerate}

第 (3) 式说明\textbf{偏向角只由两个入射角与顶角决定}，中间的折射角已被顶角关系消掉。

\begin{figure}[H]
\centering
\includegraphics[width=0.44\linewidth]{fig-3.12/fig12-6-prism.png}
\caption{棱镜主截面内的折射（讲义 p6 图 12-6）：$A$ 顶角 $\alpha$，$S\to O$ 入射，$O'\to S'$ 出射，$\delta$ 偏向角，$i,i'$ 入射角，$r,r'$ 折射角，$N,N'$ 法线。}
\end{figure}

\textbf{最小偏向角出现在 $i=i'$ 处。}把 $i'$ 看成 $i$ 的函数，逐层求导：$\cos i\dd i=n\cos r\dd r$、$\cos i'\dd i'=-n\cos r'\dd r$，两式相除消去 $\dd r$ 得 $\dd i'/\dd i=-\cos i\cos r'/(\cos i'\cos r)$。$\delta=i+i'-\alpha$ 取极值要求 $\dd\delta/\dd i=1+\dd i'/\dd i=0$，即

\[
\cos i\cos r'=\cos i'\cos r .
\]

两边平方，用 $\cos^2=1-\sin^2$ 展开，再代 $\sin i=n\sin r$、$\sin i'=n\sin r'$，化简后只剩 $(1-1/n^2)\sin^2 i=(1-1/n^2)\sin^2 i'$；因 $n\neq1$，得 $\sin i=\sin i'$，在 $0<i,i'<90^{\circ}$ 内正弦单调，故

\[
\boxed{\,i=i'\,}
\]

代回顶角关系得 $r=r'=\alpha/2$，代回偏向角得 $\delta_{\min}=2i-\alpha$，即 $i=(\delta_{\min}+\alpha)/2$。再从 $\sin i=n\sin r$：

\[
\boxed{\;
n=\frac{\sin\dfrac{\delta_{\min}+\alpha}{2}}{\sin\dfrac{\alpha}{2}}
\;}
\qquad\text{（讲义 (12-3)）}
\]

\textbf{这个驻点为什么是极小}：掠入射（$i\to90^{\circ}$）或掠出射时，光线被极度弯折，$\delta$ 迅速增大；$\delta(i)$ 在定义域内只有一个驻点，两端都大，故驻点只能是极小。

\textbf{它只要求两个角}：顶角 $\alpha$（J1 已测）与最小偏向角 $\delta_{\min}$（本环测）。\textbf{不需要知道棱镜的尺寸、厚度、材料}——这是最小偏向角法最漂亮的地方。

\callout{\textbf{为什么两个位置的读数差是 $2\delta_{\min}$。}入射光方向固定：位置 1 时出射光相对入射光偏 $+\delta_{\min}$；把棱镜转到另一面作入射面后，出射光偏 $-\delta_{\min}$。望远镜要从一侧的出射光转到另一侧，自然转过 $2\delta_{\min}$。于是
\[
\delta_{\min}=\frac12\times\frac{|\phi_2-\phi_1|+|\phi_2'-\phi_1'|}{2}
=\frac14\bigl(|\phi_2-\phi_1|+|\phi_2'-\phi_1'|\bigr).
\]
结构与 J1 完全相同（都是 $\tfrac14$ 括号里两差相加），但\textbf{物理来源不同}：J1 的 $\tfrac12$ 来自「两反射面夹角 $2\alpha$」，这里的 $\tfrac12$ 来自「偏向角一正一负」。判断方法看角标含义——同一束光的两个\textbf{反射}像用前者，同一入射光、棱镜\textbf{换面}的两个\textbf{透射}像用后者。}

\subhdr{③ 手怎么动：盯住「停顿点」}

\textbf{第 1 步：找到出射光。}转载物台到棱镜一个面朝向准直管，转望远镜找到狭缝像（这是\textbf{透射}像，比反射像暗，别和反射像混）。找到后锁 17 号，让望远镜先稳住。

\textbf{第 2 步：慢慢转载物台，跟踪像。}松 17 号，手扶望远镜始终跟着像走，同时一点一点转载物台改变入射角。像若远离入射光方向，偏向角在增大；像若靠近，偏向角在减小。朝\textbf{偏向角减小}的方向转。

\textbf{第 3 步：抓折返点。}转到某个位置，像会\textbf{停住}，然后朝相反方向移动——这个折返点就是最小偏向角的位置。

\textbf{第 4 步：锁台、微调、读数。}锁住载物台（7 号），用微调（24 号）把像对到竖直叉丝中心，读两个游标。

\textbf{第 5 步：换面重复。}把棱镜转到另一面作入射面，用同样办法再找一次折返点，再读两个游标。四个数成一组。

\begin{figure}[H]
\centering
\includegraphics[width=0.50\linewidth]{fig-3.12/fig12-13-mindev.png}
\caption{最小偏向角的两次测量（讲义 p12 图 12-13）：(a) 光从 $AB$ 入、$AC$ 出；(b) 同一入射光改从 $AC$ 入、$AB$ 出。两次出射方向各偏 $\delta_{\min}$，夹角 $2\delta_{\min}$。}
\end{figure}

\warnbox{\textbf{三个操作要点。}① \textbf{台面转得越慢越好}：越靠近极小点，像移动越慢，转快了会直接跳过折返点，读到的 $\delta$ 偏大。② \textbf{望远镜要始终跟着像}，不要转一下松手再找，那样会跟丢。③ \textbf{找的是透射像不是反射像}：透射像较暗、位置更偏，若把 $AB$ 面的反射像当成出射光，两位置夹角会完全不对。}

\subhdr{④ 当场记这些数}

\begin{center}
\begin{minipage}{\linewidth}
\small
\begin{tabular}{@{}p{1.4cm}p{2.7cm}p{2.7cm}p{2.7cm}p{2.7cm}p{2.0cm}@{}}
\toprule
\thead{次数} & \thead{$\phi_1$（位置 I）} & \thead{$\phi_1'$} & \thead{$\phi_2$（位置 II）} & \thead{$\phi_2'$} & \thead{备注} \\
\midrule
\rowstrut 1 & \blankcell & \blankcell & \blankcell & \blankcell & \\
\rowstrut 2 & \blankcell & \blankcell & \blankcell & \blankcell & \\
\rowstrut 3 & \blankcell & \blankcell & \blankcell & \blankcell & \\
\bottomrule
\end{tabular}
\end{minipage}
\end{center}

栏目与 J1 完全一样，但含义不同：这里的「位置 I / II」是\textbf{同一束入射光、棱镜换面}后的两个出射方向，不是两个反射像。

\subhdr{⑤ 当场算}

\begin{enumerate}[label=\textbf{(\arabic*)}]
\item 每行算两个差并折回（规则同 J0.5）：$|\phi_2-\phi_1|$、$|\phi_2'-\phi_1'|$。
\item 三行取平均，得 $\overline{|\Delta_1|}$、$\overline{|\Delta_2|}$。
\item $\delta_{\min}=\tfrac14\bigl(\overline{|\Delta_1|}+\overline{|\Delta_2|}\bigr)$。
\item 代 $n=\sin\frac{\delta_{\min}+\alpha}{2}\big/\sin\frac{\alpha}{2}$，其中 $\alpha$ 用 J1 的结果。
\end{enumerate}

\textbf{例算（示例数据，非实测）}：$\alpha=60^{\circ}00'$，$\phi_1=100^{\circ}00'$、$\phi_1'=280^{\circ}00'$、$\phi_2=174^{\circ}22'$、$\phi_2'=354^{\circ}21'$。

\[
|\phi_2-\phi_1|=74^{\circ}22',\qquad |\phi_2'-\phi_1'|=74^{\circ}21',
\]
\[
\delta_{\min}=\frac{74^{\circ}22'+74^{\circ}21'}{4}=37^{\circ}10.8',
\qquad
n=\frac{\sin 48^{\circ}35.4'}{\sin 30^{\circ}}=\frac{0.7500}{0.5}=1.500 .
\]

\textbf{现场数量级判断}：$n\approx1.5$ 的玻璃，$\delta_{\min}$ 应在 $37^{\circ}$--$39^{\circ}$，两位置读数差应在 $74^{\circ}$ 附近。若读数差小于 $60^{\circ}$ 或大于 $100^{\circ}$，几乎一定有错——最可能是把反射像当成了透射像，或两次不是「同一束入射光换面」。

\begin{center}
\begin{minipage}{\linewidth}
\begin{tabular}{@{}p{2.2cm}p{3.0cm}p{3.6cm}p{4.6cm}@{}}
\toprule
\thead{$n$} & \thead{$\delta_{\min}$} & \thead{读数差 $2\delta_{\min}$} & \thead{材料参考} \\
\midrule
$1.50$ & $37^{\circ}11'$ & $74^{\circ}22'$ & 冕牌玻璃量级 \\
$1.52$ & $38^{\circ}56'$ & $77^{\circ}51'$ & 普通光学玻璃 \\
$1.55$ & $41^{\circ}37'$ & $83^{\circ}13'$ & \\
$1.60$ & $46^{\circ}16'$ & $92^{\circ}31'$ & 火石玻璃量级 \\
\bottomrule
\end{tabular}
\end{minipage}
\end{center}

\textbf{$n$ 报几位}：$\alpha$ 与 $\delta_{\min}$ 只到 $0.5'$--$1'$，所以 $n$ 报\textbf{三位小数}（如 $1.512$）足够。给个量级感觉：$\delta_{\min}$ 差 $1'$（$=2.9\times10^{-4}$ rad）只让 $n$ 变 $0.0002$——\textbf{$n$ 的误差主要不是读数，而是仪器没调好}。这就是 J0 值得花五十分钟的原因。

\subhdr{⑥ 写进报告的一句话}

\callout{
「用最小偏向角法测得汞灯绿线（$546.1$ nm）在三棱镜中的最小偏向角 $\delta_{\min}=37^{\circ}10.8'$，结合顶角 $\alpha=60^{\circ}00'$ 得折射率 $n=1.500$；误差主要来自最小偏向角位置的判断（折返点附近 $\delta$ 变化平缓，位置不易精确重复）与度盘估读。」
}
\warnbox{
\textbf{台前只看这一段}：转载物台\textbf{慢慢}转、望远镜\textbf{跟着}像走，像\textbf{停住再折返}的那一点就是 $\delta_{\min}$；锁台、微调对叉丝、读两个游标；换棱镜另一面再做一次；$\delta_{\min}=\tfrac14(|\phi_2-\phi_1|+|\phi_2'-\phi_1'|)$，$n=\sin\frac{\delta_{\min}+\alpha}{2}/\sin\frac{\alpha}{2}$。
}
%=====================================================================
\clearpage
\hdr{J3}{光栅：一、二级衍射角}

\subhdr{① 这一环交出什么}

\begin{center}
\begin{minipage}{\linewidth}
\begin{tabular}{@{}p{3.6cm}p{6.0cm}p{4.4cm}@{}}
\toprule
\thead{产出} & \thead{合格判据} & \thead{交给谁} \\
\midrule
$\varphi_1,\varphi_2$（紫、绿；一、二级） & 一级 $7.5^{\circ}$（紫）、$9.4^{\circ}$（绿）；二级 $15.2^{\circ}$、$19.1^{\circ}$ & J4 算 $d,N,D$ \\
\bottomrule
\end{tabular}
\end{minipage}
\end{center}

用时约 65 分钟：调光栅 20 分钟，八个方位读数 45 分钟。

\subhdr{② 为什么这样能测：从双缝干涉接过来}

这一节假设你\textbf{没有学过衍射}。不需要——光栅方程可以从你在第五章学过的双缝干涉直接推出来，只多一步。

\subhdr{2.1 先把「多走了多少路」算出来}

两条缝相距 $d$。平行光垂直照到缝上，两缝的光往同一个方向走——这个方向与光栅法线的夹角记作 $\varphi$。

本来两束光是「同时出发」的；斜着走的时候，靠后的那条缝射出的光要多走一段路，这段路就是\textbf{光程差}：

\[
\Delta=d\sin\varphi .
\]

这就是全部的几何：从一条缝的光线上，向另一条光线作一条垂线，那个直角三角形的斜边是 $d$、对边就是 $d\sin\varphi$。

\subhdr{2.2 什么时候亮：光程差是波长的整数倍}

你在干涉那一章学过的判据是：两束光的光程差是波长的整数倍时，波峰对波峰，\textbf{加强}（亮）；是半波长的奇数倍时，波峰对波谷，\textbf{相消}（暗）。于是亮纹条件写成

\[
\boxed{\,d\sin\varphi=k\lambda\,}\qquad k=0,\pm1,\pm2,\dots
\]

这就是\textbf{光栅方程}。$k$ 叫\textbf{衍射级次}：$k=0$ 是零级（$\varphi=0$，各波长都不偏折，全落在中央一条上），$k=\pm1,\pm2$ 对称分布在两侧。

\textbf{光栅与双缝只差「缝的条数」。}双缝的两条缝也能产生这个条件，但亮纹很宽、很暗；光栅有成千条缝，同一方向上所有缝的光都同相叠加，于是亮纹变得\textbf{又亮又细}。所以光栅能分辨很靠近的两条谱线（钠双线 $589.0$ 与 $589.6$ nm 就靠这个分开）。这一点口试会问，记住「多缝同相叠加 $\Rightarrow$ 又亮又细」即可，不需要算强度分布。

\begin{figure}[H]
\centering
\includegraphics[width=0.72\linewidth]{fig-3.12/fig12-7-grating.png}
\caption{光栅衍射的几何（讲义 p7 图 12-7）：间距 $d$ 的相邻缝，沿与法线成 $\varphi$ 的方向，光程差为 $d\sin\varphi$；对称出现零级与 $\pm1,\pm2$ 各级主极大。}
\end{figure}

\subhdr{2.3 为什么紫在内、绿在外}

同一块光栅、同一级次，$\lambda$ 越小则 $\sin\varphi$ 越小，也就是\textbf{偏折越小、越靠近零级}。汞灯里紫线 $\lambda=435.8$ nm、绿线 $546.1$ nm，所以紫在内、绿在外。这和棱镜正好相反（棱镜靠折射率随波长变，紫的 $n$ 大、偏得多，紫在外）——口试的第 11 问就是这个对比。

\textbf{还没测就能知道答案}：若光栅是 $300$ 线/mm，$d=3333$ nm，则

\begin{center}
\begin{minipage}{\linewidth}
\begin{tabular}{@{}p{2.6cm}p{3.4cm}p{3.4cm}p{3.4cm}@{}}
\toprule
\thead{谱线} & \thead{$\sin\varphi_1=\lambda/d$} & \thead{一级 $\varphi_1$} & \thead{二级 $\varphi_2$} \\
\midrule
紫 $435.8$ nm & $0.1307$ & $7.51^{\circ}$ & $15.16^{\circ}$ \\
绿 $546.1$ nm & $0.1638$ & $9.43^{\circ}$ & $19.13^{\circ}$ \\
\bottomrule
\end{tabular}
\end{minipage}
\end{center}

现场拿这张表对：一级绿线若出现在 $3^{\circ}$ 出头，你的光栅是 $100$ 线/mm；若出现在 $19^{\circ}$ 附近，是 $600$ 线/mm。\textbf{先看光栅框上标的线数}，再对表。

\subhdr{③ 手怎么动}

\textbf{第 1 步：对零级。}把望远镜转到正对准直管，让竖直叉丝与狭缝像重合，锁 17 号。这一步定的是「零级方向」。

\textbf{第 2 步：放光栅。}取下棱镜，把光栅放上载物台：从上方看，$b_1b_2$ 的连线应被栅面\textbf{垂直平分}。

\textbf{第 3 步：调栅面垂直光轴。}慢慢转载物台让栅面朝向望远镜，找栅面反射的绿色「$+$」像；调 \textbf{$b_1$、$b_2$}（\textbf{不要调 $b_3$}）让绿像与上叉丝重合。此时栅面垂直望远镜光轴。

\textbf{第 4 步：零级居中。}调 27 号（准直管光轴竖直），让零级亮线被横丝\textbf{垂直平分}，且像最亮最清晰。

\textbf{第 5 步：调等高。}松 17 号，转望远镜看 $\pm1,\pm2$ 各级谱线是否\textbf{等高}；不等高就调 \textbf{$b_3$}，直到各级在同一高度——这一步让栅缝平行主轴。

\begin{figure}[H]
\centering
\includegraphics[width=0.36\linewidth]{fig-3.12/fig12-14-gratingadj.png}
\caption{光栅在载物台上的方位与三个螺钉（讲义 p13 图 12-14）：$b_1,b_2$ 调栅面垂直光轴，$b_3$ 调栅缝方向；$b_1b_2$ 连线被栅面垂直平分。}
\end{figure}

\callout{\textbf{讲义 p13 问了一句：「只需调光栅的一侧，为什么不用转 $180^{\circ}$？」}因为光栅只有\textbf{一个面}工作，测量时也不会翻面；这里只有一个未知量（栅面与轴的夹角），一面就够。而平面镜那一步之所以要转 $180^{\circ}$，是因为那里要同时定出「台面倾斜」与「望远镜轴倾斜」两个未知量，必须用两个面的同向/反向信息把它们分开。栅面垂直望远镜轴 $\Leftrightarrow$ 栅面平行主轴，衍射角读出来才是真的。}

\subhdr{第 6 步：测 $\pm k$ 级}

转动望远镜，把竖直叉丝依次对准 $k=\pm1$、$\pm2$ 各条谱线的中心，\textbf{每条读两个游标}。一个方位四个动作：对准、锁 17、读两游标、记下来。

顺序建议：先测绿线（亮、好找）的 $+1$、$-1$、$+2$、$-2$，再测紫线同样四个。紫线暗，若一时找不到，先把缝宽略微调大一点点再找，找到后调回去。

\begin{figure}[htbp]
\centering
\includegraphics[width=0.62\linewidth]{fig-3.12/fig12-9-spectra.png}
\caption{汞灯的多级光谱（讲义 p8 图 12-9）：中央是零级（白亮），两侧对称出现一级、二级。级次越高越暗、越散开，所以二级要仔细找。}
\end{figure}

\subhdr{④ 当场记这些数}

\begin{center}
\begin{minipage}{\linewidth}
\scriptsize\setlength{\tabcolsep}{3pt}
\begin{tabular}{@{}p{2.2cm}p{1.55cm}p{1.55cm}p{1.55cm}p{1.55cm}p{1.55cm}p{1.55cm}p{1.55cm}p{1.55cm}@{}}
\toprule
\thead{颜色} & \thead{$\theta_1$} & \thead{$\theta_1'$} & \thead{$\theta_2$} & \thead{$\theta_2'$} & \thead{$\theta_{-1}$} & \thead{$\theta_{-1}'$} & \thead{$\theta_{-2}$} & \thead{$\theta_{-2}'$} \\
\midrule
\rowstrut 紫 $435.8$ nm & \blankcell & \blankcell & \blankcell & \blankcell & \blankcell & \blankcell & \blankcell & \blankcell \\
\rowstrut 绿 $546.1$ nm & \blankcell & \blankcell & \blankcell & \blankcell & \blankcell & \blankcell & \blankcell & \blankcell \\
\bottomrule
\end{tabular}
\end{minipage}
\end{center}

\textbf{栏目含义}：$\theta_k$ 与 $\theta_k'$ 是\textbf{同一级}谱线在两个游标上的读数；下标带负号的是零级\textbf{另一侧}的同级谱线。一行 8 个数，两行共 16 个数——这就是手册表 4 的全部内容。

\subhdr{⑤ 当场算}

\[
\varphi_k=\frac14\bigl(|\theta_k-\theta_{-k}|+|\theta_k'-\theta_{-k}'|\bigr).
\]

\textbf{为什么是 $\tfrac14$}：$\pm k$ 两级的夹角是 $2\varphi_k$（对称分布），这是第一个 $\tfrac12$；两个游标各给一次、取平均，是第二个 $\tfrac12$。和 J1、J2 同一个结构。

\textbf{为什么用 $\pm k$ 平均而不是各自算}：零级在度盘上的读数没有人去测，它可能整体偏了一个小角 $\beta$（栅面没完全垂直轴，或望远镜零位没对准）。设偏置为 $\beta$，则 $+k$ 级偏 $+\beta$、$-k$ 级偏 $-\beta$；两者相减，$\beta$ 自动抵消——与双游标消偏心是同一个思路：\textbf{用对称的两个测量把共同的偏置减掉}。

\textbf{例算}：测得绿线 $\theta_1=9^{\circ}26'$、$\theta_{-1}=350^{\circ}34'$（零级偏置已包含在这两个数里），则

\[
|\theta_1-\theta_{-1}|=|9^{\circ}26'-350^{\circ}34'|=341^{\circ}08'
\ \xrightarrow{\ 360^{\circ}-\ \ }18^{\circ}52',
\]

这个 $18^{\circ}52'$ 是 $2\varphi_1$，故 $\varphi_1=9^{\circ}26'$——与预期 $9.43^{\circ}$ 相符。（真实数据里两个游标各给一次，四项一起平均。）

\textbf{最灵的自检}：光栅方程要求 $\sin\varphi_k=k\sin\varphi_1$，所以

\[
\frac{\sin\varphi_2}{\sin\varphi_1}\approx2 .
\]

算出 $1.8$ 或 $2.2$，说明某一级读错了游标或认错了谱线（最常见：把二级的某一色当成了一级）。这条自检比任何复核都管用，算完立刻做。

\subhdr{⑥ 写进报告的一句话}

\callout{
「以汞灯为光源，测得一级衍射角 $\varphi_1(\text{紫})=7.5^{\circ}$、$\varphi_1(\text{绿})=9.4^{\circ}$，二级 $\varphi_2(\text{紫})=15.2^{\circ}$、$\varphi_2(\text{绿})=19.1^{\circ}$；$\sin\varphi_2/\sin\varphi_1\approx2.00$，说明两级读数自洽。」
}
\warnbox{
\textbf{台前只看这一段}：光栅放好后 $b_1b_2$ 调「栅面垂直光轴」、$b_3$ 调「各级等高」；先对零级，再测 $\pm1$、$\pm2$，每条读两个游标；$\varphi_k=\tfrac14\bigl(|\theta_k-\theta_{-k}|+|\theta_k'-\theta_{-k}'|\bigr)$；算完立刻查 $\sin\varphi_2/\sin\varphi_1\approx2$。
}
%=====================================================================
\clearpage
\hdr{J4}{光栅常量、线数与角色散}

\subhdr{① 这一环交出什么}

\begin{center}
\begin{minipage}{\linewidth}
\begin{tabular}{@{}p{3.6cm}p{4.2cm}p{6.2cm}@{}}
\toprule
\thead{产出} & \thead{合格判据} & \thead{数据来源} \\
\midrule
光栅常量 $d$、线数 $N$、角色散 $D$ & $N$ 接近整数（$100/300/600$）；一级 $D\approx1'/\mathrm{nm}$ & 全部用 J3 的 $\varphi_k$，\textbf{不需要新测量} \\
\bottomrule
\end{tabular}
\end{minipage}
\end{center}

用时约 30 分钟，全部是计算与核对。这一环可以离开仪器做，但\textbf{如果发现 \S ⑤ 的自检不过，趁仪器还在赶紧回去补测}。

\subhdr{② 为什么这样能算}

\textbf{光栅常量}：光栅方程里已知 $\lambda$、测得 $\varphi_1$，直接解出

\[
d=\frac{\lambda}{\sin\varphi_1}.
\]

\textbf{线数}：出厂标的是「每毫米多少条缝」$N$，而公式里用的是缝距 $d$。两者的关系是 $d=1/N$（$d$ 以 mm 计），也就是

\[
N=\frac{1}{d}\quad(d\ \text{以 mm 计}),\qquad
N=\frac{1}{1000\,d}\quad(d\ \text{以 m 计}).
\]

\textbf{角色散}：定义为「单位波长差对应的角间距」，

\[
D=\frac{\Delta\varphi}{\Delta\lambda}
\quad\text{（直接测量法：用一级的紫线与绿线）},
\]
\[
D=\frac{\dd\varphi}{\dd\lambda}=\frac{k}{d\cos\varphi}
\quad\text{（由光栅方程求导：$d\sin\varphi=k\lambda$ 两边对 $\lambda$ 求导）}.
\]

两个式子应互相印证。$300$ 线/mm 的光栅，一级角色散约 $1'/\mathrm{nm}$——记住这个量级，它是判断对错的锚。

\subhdr{③ 不需要新操作}

这一环的全部输入是 J3 的八个方位读数。唯一要确认的是：\textbf{八个读数对应同一块光栅、同一套调节}；中途若动过 16 号或重新调过 $b_1b_2$，前后数据不能混着用。

\subhdr{④ 汇总表（把 J3 的结果抄进来算）}

\begin{center}
\begin{minipage}{\linewidth}
\begin{tabular}{@{}p{3.4cm}p{3.4cm}p{3.4cm}p{4.4cm}@{}}
\toprule
\thead{量} & \thead{公式} & \thead{测得值} & \thead{与预期是否相符} \\
\midrule
\rowstrut $d$（一级绿线） & $\lambda/\sin\varphi_1$ & & $3333$ nm（$300$ 线/mm） \\
\rowstrut $d'$（二级绿线核对） & $2\lambda/\sin\varphi_2$ & & 与 $d$ 差 $<1\%$ \\
\rowstrut $N$ & $1/d$ & & $100/300/600$ \\
\rowstrut $D$（一级紫绿） & $\Delta\varphi/\Delta\lambda$ & & $\approx1.0\ '/\mathrm{nm}$ \\
\rowstrut $D$（解析式核对） & $k/(d\cos\varphi)$ & & 与上一行差 $<1\%$ \\
\bottomrule
\end{tabular}
\end{minipage}
\end{center}

\subhdr{⑤ 当场算：一步一步}

\textbf{第 1 步：$d$ 用一级绿线。}手册明确指定一级绿线（$\lambda=546.1$ nm）：

\[
d=\frac{546.1\ \mathrm{nm}}{\sin 9^{\circ}26'}=\frac{546.1}{0.16387}=3333\ \mathrm{nm}=3.333\ \mu\mathrm{m}.
\]

\textbf{第 2 步：用二级做独立核对。}

\[
d'=\frac{2\lambda}{\sin\varphi_2}=\frac{1092.2\ \mathrm{nm}}{\sin 19^{\circ}08'}=\frac{1092.2}{0.3277}=3333\ \mathrm{nm}\ \checkmark
\]

理论上 $d'=d$。实测应相差在 $1\%$ 以内；差很多说明某一级角度测偏了（多半是栅面没垂直光轴，或谱线不等高）。

\textbf{第 3 步：换成线数并与出厂值对照。}

\[
N=\frac{1}{d}=\frac{1}{3.333\times10^{-3}\ \mathrm{mm}}=300\ \text{线/mm}.
\]

$N$ 应接近整数。算出 $293$ 或 $307$ 是正常的测量偏差；算出 $2930$ 说明单位换算错了 10 倍。

\textbf{第 4 步：角色散用一级紫绿两线。}

\[
D=\frac{\varphi_{1,\text{绿}}-\varphi_{1,\text{紫}}}{546.1-435.8}
=\frac{9^{\circ}26'-7^{\circ}31'}{110.3\ \mathrm{nm}}
=\frac{115'}{110.3\ \mathrm{nm}}=1.04\ '/\mathrm{nm}.
\]

与解析式 $D=\dfrac{1}{d\cos 9^{\circ}26'}=\dfrac{1}{3333\times0.9865}=1.045'/\mathrm{nm}$ 一致 $\checkmark$。

\warnbox{\textbf{三处最容易出错的地方}：① 把 $N$（线/mm）当成 $d$ 直接代进 $d\sin\varphi=k\lambda$；② $d$ 用 $\mu$m、$\lambda$ 用 nm，忘了换算（$1\ \mu\mathrm{m}=1000$ nm）；③ 漏掉 $N=1/(1000d)$ 里的 $1000$，算出的 $d$ 差一千倍。统一做法：\textbf{先把 $d$ 换到 nm}，再与 $\lambda$ 一起代。}

\begin{center}
\begin{minipage}{\linewidth}
\begin{tabular}{@{}p{2.2cm}p{3.0cm}p{3.4cm}p{4.6cm}@{}}
\toprule
\thead{$N$（线/mm）} & \thead{$d=1/N$} & \thead{一级绿线 $\varphi_1$} & \thead{量级感觉} \\
\midrule
$100$ & $10.000\ \mu\mathrm{m}$ & $3.13^{\circ}$ & 一级很近，容易与零级混 \\
$300$ & $3.333\ \mu\mathrm{m}$ & $9.43^{\circ}$ & 本实验最常见 \\
$600$ & $1.667\ \mu\mathrm{m}$ & $19.13^{\circ}$ & 一级已经很散开 \\
\bottomrule
\end{tabular}
\end{minipage}
\end{center}

\textbf{光栅光谱与棱镜光谱的区别}（口试第 11 问）：

\begin{center}
\begin{minipage}{\linewidth}
\begin{tabular}{@{}p{2.4cm}p{5.4cm}p{5.4cm}@{}}
\toprule
\thead{比较项} & \thead{光栅} & \thead{棱镜} \\
\midrule
分光机制 & 多缝干涉（$d\sin\varphi=k\lambda$） & 折射率随波长变（色散） \\
谱线排列 & 长波偏折大（\textbf{红在外}），近似等间距 & 短波偏折大（\textbf{紫在外}），不等间距 \\
级次 & 有多级，高级次会重叠 & 只有一套光谱 \\
\bottomrule
\end{tabular}
\end{minipage}
\end{center}

\subhdr{⑥ 写进报告的一句话}

\callout{
「由一级绿线得光栅常量 $d=3.333\ \mu\mathrm{m}$，对应线数 $N=300$ 线/mm；用二级复核得 $d'=3.333\ \mu\mathrm{m}$，两者相差小于 $1\%$；由一级紫绿两线得角色散 $D=1.04'/\mathrm{nm}$，与解析式 $k/(d\cos\varphi)$ 一致。」
}
\warnbox{
\textbf{台前只看这一段}：$d=\lambda/\sin\varphi_1$（用一级\textbf{绿}线，$\lambda=546.1$ nm）；$d'=2\lambda/\sin\varphi_2$ 核对；$N=1/d$（$d$ 换 mm）；$D=\dfrac{\varphi_{\text{绿}}-\varphi_{\text{紫}}}{110.3\ \mathrm{nm}}$；先把 $d$ 换到 nm 再代。
}
%=====================================================================
\clearpage
\hdr{A}{附录 A：想深一层（不读也能做完实验）}

这三节回答「为什么这个仪器这样做」，与操作无关。现场时间紧就跳过，报告写完再看。

\subhdr{A.1 图 12-1 的 28 个编号：你真正会动的只有 12 个}

不必全记。按「什么时候动它」分三组，编号与原名不变：

\begin{center}
\begin{minipage}{\linewidth}
\small
\begin{tabular}{@{}p{0.6cm}p{6.0cm}p{7.4cm}@{}}
\toprule
\thead{编号} & \thead{名称（英文原名）} & \thead{什么时候动它} \\
\midrule
\multicolumn{3}{@{}l}{\textbf{A. 望远镜这一路}}\\
8 & 望远镜 \emph{telescope} & 转它去找像、找谱线 \\
9 & 目镜筒锁紧 \emph{locking screw of eyepiece} & 松开后目镜筒才能前后移 \\
11 & 目镜视度调节手轮 \emph{adjustment hand wheel} & \textbf{最先动它}：把叉丝本身看清 \\
12 & 望远镜光轴竖直调节 \emph{vertical adjustment screw} & \textbf{各半法调它} \\
13 & 望远镜光轴水平调节 \emph{horizontal adjustment screw} & 目测粗调时对准准直管 \\
15 & 望远镜微调 \emph{fine tuning screw} & 把像对到叉丝中心 \\
17 & 望远镜锁紧 \emph{locking screw of telescope} & 对准后锁死 \\
\midrule
\multicolumn{3}{@{}l}{\textbf{B. 准直管这一路}}\\
1 & 狭缝装置 \emph{slit device} & 转它让狭缝像竖直 \\
2 & 狭缝锁紧 \emph{slit device locking screw} & 松—移—紧一次完成，之后别碰 \\
26 & 准直管光轴水平调节 \emph{horizontal screw} & 目测粗调时对准望远镜 \\
27 & 准直管光轴竖直调节 \emph{vertical screw} & 步骤 5 调它：像落在视场中央 \\
28 & 缝宽调节手轮 \emph{hand wheel of slit width} & 调到约 $0.5$ mm（范围 $0.2$--$2$ mm） \\
\midrule
\multicolumn{3}{@{}l}{\textbf{C. 载物台、度盘与游标}}\\
5 & 载物台 \emph{specimen stage} & 放平面镜、三棱镜、光栅 \\
6 & 载物台调平螺钉 \emph{leveling bolt}（三个 $a_1a_2a_3$） & \textbf{各半法调它}；三螺钉成正三角形 \\
16 & 游标盘与度盘的锁紧 \emph{locking screw}（背面） & \textbf{决定「谁跟谁一起转」}；一组测量开始后不许改 \\
21 & 游标 \emph{vernier}（两个，相差 $180^{\circ}$） & \textbf{读数看它} \\
22 & 度盘 \emph{dial plate}（$0$--$360^{\circ}$） & 读数看它 \\
24 & 游标盘微调 \emph{fine tuning screw of alidade} & 把像对到叉丝中心 \\
25 & 游标盘锁紧 \emph{locking screw of alidade} & 对准后锁死 \\
7 & 载物台锁紧 & 找到最小偏向角后锁台 \\
\bottomrule
\end{tabular}
\end{minipage}
\end{center}

%=====================================================================
\needspace{16\baselineskip}\hdr{B}{附录 B：两份材料说法不一致时，以这条为准}

讲义（15 页）与预习（3 页）有八处彼此不一致或没写清楚。现场遇到就照下表，不要在仪器前纠结。

\begin{center}
\begin{minipage}{\linewidth}
\small
\begin{tabular}{@{}p{4.0cm}p{5.2cm}p{5.0cm}@{}}
\toprule
\thead{问题} & \thead{材料怎么说} & \thead{本讲口径} \\
\midrule
重复几次 & 正文说 five times，预习表只留 3 行 & 按 3 组做；老师要 5 次就加两行 \\
光栅线数 & 只给 $N=1/(1000d)$，没给具体数 & 先看光栅框上的标注；未标注按 $300$ 线/mm 预估角度 \\
棱镜是否等边 & 未写 & 按等边（$\alpha=60^{\circ}$）做例算与判据 \\
游标「取零线右边」 & 与它自己的算例矛盾 & 以算例为准：取零线\textbf{已越过}的那条 \\
$(12\text{-}1)(12\text{-}2)$ & 按 $\angle1=\angle2$ 代入会得 $\Phi=\theta_1$，与结论矛盾 & 属笔误；结论 $\phi=(\theta_1+\theta_2)/2$ 不变（附录 A.2 重推） \\
\emph{dispersive power} & 引言用这个词，正文定义的却是 $D=\Delta\varphi/\Delta\lambda$ & 一律按\textbf{角色散}处理 \\
图 12-12 角标 & 第二个游标印成 $\phi_2''$ & 统一用 $\phi_2'$ \\
两套螺钉编号 & 图 12-10 用 $a,b,c$，图 12-11 用 $a_1,a_2,a_3$ & 不是同一套；现场先看载物台刻字，按「只改一个面」的原则挑 \\
钠灯 & 器材清单列了，流程全用汞灯 & 本讲不给钠灯操作，不编造步骤 \\
\bottomrule
\end{tabular}
\end{minipage}
\end{center}

%=====================================================================
\needspace{16\baselineskip}\hdr{C}{附录 C：最贵的十个错}

\begin{center}
\begin{minipage}{\linewidth}
\small
\begin{tabular}{@{}p{0.6cm}p{5.6cm}p{7.6cm}@{}}
\toprule
\thead{\#} & \thead{错法} & \thead{后果与正确做法} \\
\midrule
1 & 只读一个游标 & 引入最多 $1'$ 的偏心误差。\textbf{每读一次方位必须两个游标都读} \\
2 & 「像找到了」当成「调好了」 & 绿十字模糊或带视差就开始测，所有角量带固定偏差。判据是\textbf{清晰 + 无视差 + 与上叉丝重合}三件齐 \\
3 & 各半法一次调到位 & 一面重合、另一面偏差更大且不收敛。正确是\textbf{先减半、再归零} \\
4 & 折回 $360^{\circ}$ 只在平均后做一次 & 某行原始差 $240^{\circ}$ 当 $240^{\circ}$ 参与平均，$\alpha$ 错几十度。\textbf{每行都判断} \\
5 & $d$ 与 $N$ 混用 / 单位没换 & 结果差 $10^3$ 量级。\textbf{先把 $d$ 换到 nm} \\
6 & 调 $a_2$ 想「一起摆平」 & 两个面被同时扰动。\textbf{只调 $a_1,a_3$} \\
7 & 测最小偏向角时转太快 & 跳过折返点，$\delta$ 偏大。\textbf{慢转、盯住停顿} \\
8 & 光栅只调 $b_1,b_2$ 没调 $b_3$ & 谱线不等高，叉丝与斜谱线交点含糊 \\
9 & 中途动了 16 号螺钉 & 前后读数不属同一度盘，整组作废。\textbf{一组开始后别碰} \\
10 & 读数写成 $335^{\circ}47'$ 而实际 $334^{\circ}47'$ & 没做「结果落在两刻线之间」的自检 \\
\bottomrule
\end{tabular}
\end{minipage}
\end{center}

\warnbox{\textbf{时间不够时只保三件事}：① 望远镜聚焦到无视差；② 望远镜轴垂直主轴（镜面两面都重合）；③ 每次读数都读两个游标。这三件保证数据的系统性；其余的（少测一组、二级没测）只是精度损失，可以事后说明。反过来，这三件缺任何一件，数据不可救。}

\subhdr{A.2 自准直为什么比肉眼准：灵敏度放大一倍}

设平面镜相对光轴偏了 $\gamma$，由反射定律，返回光的方向偏了 $2\gamma$，像点在上叉丝旁移开

\[
s=2f\gamma .
\]

取 $f=200$ mm、$\gamma=1'=2.91\times10^{-4}$ rad，则 $s=2\times200\times2.91\times10^{-4}=0.116$ mm——肉眼清楚可见。

这就是自准直的价值：\textbf{角度的判据被放大了一倍，再被焦距换成一段可以直接看的位移}。$1'$ 的偏角在视场里是一段看得见的距离，所以「像与上叉丝重合」比「看上去大致对准」精确得多。反过来说，$f$ 越大越灵敏，但视场越小、越难找像——这就是为什么粗调必须先把像找进视场。

\subhdr{A.3 偏心差的一阶推导（为什么两个游标取平均就消掉）}

设主轴中心 $O$ 在原点，度盘中心 $O'$ 在 $(e\cos\psi,\ e\sin\psi)$，度盘半径 $R$。游标 $V$ 随望远镜绕 $O$ 转到真实方位角 $\phi$，而度盘刻线从 $O'$ 读出角 $\theta$。于是

\[
V=\bigl(e\cos\psi+R\cos\theta,\ e\sin\psi+R\sin\theta\bigr).
\]

从 $O$ 看 $V$ 的方向满足

\[
\tan\phi=\frac{R\sin\theta+e\sin\psi}{R\cos\theta+e\cos\psi}
=\tan\theta+\frac{(e/R)\sin(\psi-\theta)}{\cos^2\theta}+\cdots
\]

一阶近似给出 $\phi-\theta\approx(e/R)\sin(\psi-\theta)$，即读数误差

\[
\varepsilon(\theta)=\theta-\phi\approx\frac{e}{R}\sin(\theta-\psi).
\]

两个游标装在一条直径的两端，方位差 $180^{\circ}$：

\[
\varepsilon(\theta+180^{\circ})=\frac{e}{R}\sin(\theta+180^{\circ}-\psi)=-\frac{e}{R}\sin(\theta-\psi)=-\varepsilon(\theta).
\]

\textbf{等大反号}，取平均即消：

\[
\theta_1=\phi+\varepsilon,\quad\theta_2=\phi-\varepsilon
\quad\Longrightarrow\quad
\phi=\frac{\theta_1+\theta_2}{2}.
\]

副产品：$\varepsilon=(\theta_1-\theta_2)/2$，也就是\textbf{两个游标的读数差本身量出了偏心误差}——现场可用它做自检（差应在几个角分内）。

\textbf{量级}：$e=0.02$ mm、$R=60$ mm 时，$\varepsilon_{\max}=\arctan(0.02/60)=1.15'\approx1'$，比读数精度（$0.5'$）还大；平均后只剩 $(e/R)^2$ 量级，约 $0.01''$，完全可忽略。


%=====================================================================
\needspace{16\baselineskip}\hdr{D}{附录 D：报告速查}

\subhdr{D.1 报告要交什么}

\begin{center}
\begin{minipage}{\linewidth}
\begin{tabular}{@{}p{3.6cm}p{3.2cm}p{7.0cm}@{}}
\toprule
\thead{项} & \thead{什么时候} & \thead{内容} \\
\midrule
预习 8 题 & 进实验室前 & 英文答案写进报告（下面 D.2 可直接抄） \\
数据表 & 课上 & 打印 Preview p3 或本书 J1/J2/J3 的表 \\
数据处理 & 课后 & 六项：$\alpha$、$\delta_{\min}$、$n$、$\varphi_k$、$d$、$D$（J1--J4 的 \S ⑤） \\
口试 & 课上可能被问 & 12 问要点（D.3），不必写进报告 \\
选做 & 自愿 & 3 道扩展题（D.4） \\
\bottomrule
\end{tabular}
\end{minipage}
\end{center}

\subhdr{D.2 预习 8 题：中英答案}

每题先给\textbf{可直接抄进报告的英文}（压到 50--60 词），再给\textbf{中文解释}（理解用，不必抄）。

\subhdr{Q1\quad What are the requirements for fine adjustment of the spectrometer?}

\textbf{EN.} Three conditions: (1) the collimator emits a parallel beam, i.e. its slit lies in the focal plane of the collimator lens; (2) the telescope receives parallel light, i.e. it is focused at infinity with the reticle in the focal plane of the objective; (3) both optical axes are perpendicular to the rotation axis, and the optical surface of the element under test is parallel to it.

\textbf{中文.} 三条：准直管出平行光、望远镜收平行光、两轴都垂直主轴且待测元件的光学面平行主轴。第三条最容易只做一半——前两条只保证「有平行光」，第三条才保证「角度读数有意义」。

\subhdr{Q2\quad What is the purpose of visual coarse adjustment? How to carry it out?}

\textbf{EN.} Purpose: to bring both axes roughly perpendicular to the rotation axis and the stage roughly horizontal, so that the autocollimated image can enter the small field of view; the fine screws move only a few degrees. Method: by eye set the telescope (No.\,12) and the collimator (No.\,27) at the same height facing each other, then level the stage with its three screws.

\textbf{中文.} 目的是把三条轴摆进「细调能起作用」的范围。讲义原文说粗调「是细调成功的关键」，不是客套。

\subhdr{Q3\quad Explain the principle of focusing on infinity by the self-aligning method.}

\textbf{EN.} The illuminated transparent cross acts as the object. If the reticle lies in the focal plane, its light leaves the objective parallel and, after reflection by a plane mirror, returns parallel and is focused back onto the reticle plane; image and crosshair then show no parallax. Eliminating parallax therefore focuses the telescope at infinity.

\textbf{中文.} 关键是一条链：物在焦面 $\Leftrightarrow$ 出射平行 $\Leftrightarrow$ 返回后仍聚焦在焦面 $\Leftrightarrow$ 无视差。把这条链说完整，比只背「调到这里没有视差」得分高。

\subhdr{Q4\quad Draw a sketch of the plane mirror placement; how to judge the axis is perpendicular?}

\textbf{EN.} (Sketch: the stage as a circle with three screws; the mirror standing vertically, one face toward the telescope.) Judgment: the green cross from the front face coincides with the upper crosshair; after turning the stage by $180^{\circ}$ the image from the rear face must also coincide. Both coinciding means the telescope axis is perpendicular to the rotation axis.

\textbf{中文.} 画图要素：三个螺钉、镜子竖直、一个反射面朝望远镜。判据是\textbf{两面都重合}，不是一面。

\subhdr{Q5\quad How to judge that the parallel light from the collimator is perpendicular to the rotation axis?}

\textbf{EN.} Align the telescope with the collimator and watch the slit image. Turn the slit sleeve until the image is parallel to the horizontal crosshair, then adjust screw No.\,27 until the image is centred vertically; the collimator axis is then perpendicular to the rotation axis. Finally turn the slit by $90^{\circ}$ and lock it.

\textbf{中文.} 判据两层：像\textbf{居中}（轴垂直主轴）与像\textbf{竖直}（狭缝本身竖直，这是后面读衍射角的前提）。

\subhdr{Q6\quad Sketch the reflection method for the vertex angle, and give the formula.}

\textbf{EN.} (Sketch: telescope at I for the beam from face $AB$, at II for the beam from $AC$.) Rotating a mirror by $\gamma$ turns the reflected ray by $2\gamma$; the two faces differ by $\alpha$, so the two beams differ by $2\alpha$ and the telescope turns through $2\alpha$. Hence
\[
\alpha=\frac14\bigl(|\phi_2-\phi_1|+|\phi_2'-\phi_1'|\bigr),
\]
with $\phi_1,\phi_1'$ and $\phi_2,\phi_2'$ the vernier readings at I and II.

\textbf{中文.} 图要画出准直管、顶点、两个反射面、两束反射光、望远镜两个位置与四个角标。$\tfrac14$ 来自「双面镜的 $\tfrac12$」乘「双游标的 $\tfrac12$」。

\subhdr{Q7\quad Sketch the method of measuring the minimum deflection angle.}

\textbf{EN.} (Sketch: the prism in both orientations, the two emergent directions on either side of the fixed incident direction, $2\delta_{\min}$ apart.) Rotate the stage slowly: the slit image moves, stops, then reverses; the stationary point is the minimum deviation. Read both verniers, turn the prism to the other face and repeat. The difference is $2\delta_{\min}$, so
\[
\delta_{\min}=\frac14\bigl(|\phi_2-\phi_1|+|\phi_2'-\phi_1'|\bigr).
\]

\textbf{中文.} 图要能看出「入射方向不变、出射方向在两侧各偏 $\delta_{\min}$」。文字部分必须写「像停住后折返」这个判据——那是操作的核心。

\subhdr{Q8\quad Explain the principle of measuring grating constants; give the grating equation.}

\textbf{EN.} For normal incidence, light from two adjacent slits has a path difference $d\sin\varphi$ in the direction $\varphi$; a bright line appears when this equals an integer number of wavelengths:
\[
d\sin\varphi=k\lambda\qquad(k=0,\pm1,\pm2,\dots).
\]
Measuring $\varphi_1$ for a known $\lambda$ gives $d=\lambda/\sin\varphi_1$, and the line density is $N=1/d$ ($d$ in mm).

\textbf{中文.} 答题必须把「光程差 $d\sin\varphi$」这一步写出来，只写公式不写来源通常扣分。

\needspace{20\baselineskip}\subhdr{D.3 口试 12 问：要点}

\begin{center}
\begin{minipage}{\linewidth}
\small
\begin{tabular}{@{}p{0.5cm}p{4.8cm}p{8.6cm}@{}}
\toprule
\thead{\#} & \thead{问题} & \thead{要点} \\
\midrule
1 & 分光计由几部分组成？各起什么作用？ & 五部分：望远镜（收平行光、自准直定位方向）、准直管（出平行光，把缝的位置换成方向）、载物台（放待测件、绕主轴转）、度盘与游标（把方位读成 $1'$）、底座（支撑与定位主轴） \\
2 & 主截面是什么？ & 垂直于棱镜两折射面交线的截面。只有主截面内的光满足 $\alpha=r+r'$ 与 $\delta=i+i'-\alpha$；主截面垂直主轴是全部角度公式的前提 \\
3 & 为什么要求望远镜光轴平行于主截面？怎么调？ & 因为公式在主截面内推。调法是让主截面垂直主轴：$AB$、$AC$ 两面自准直像都与上叉丝重合，只调 $a_1,a_3$ \\
4 & 自准直的步骤与「能接收平行光」的判据？ & 四步：看清叉丝 → 找绿十字 → 移目镜筒至清晰 → 消视差。判据「清晰且无视差」，因为无视差 $\Leftrightarrow$ 分划板在焦面 $\Leftrightarrow$ 出射与返回都平行 \\
5 & 为什么绿十字与\textbf{上叉丝}重合？ & 透光「十」在轴下 $a$、上叉丝在轴上 $a$，两者对称；出射平行光与轴成 $\theta=a/f$，反射只反转轴向分量，返回光向另一侧偏，落在 $+a$ \\
6 & 细调的要求有哪些？ & 同 Q1 三条：准直管出平行光、望远镜聚焦无穷远、两轴垂直主轴 \\
7 & 找不到平面镜的「$+$」像怎么办？ & 回目测粗调：查电源与绿 LED、镜子是否放平且反射面朝望远镜、12/13 号是否偏太多、9 号镜筒是否把像推出视场 \\
8 & 测顶角时找不到棱镜反射像怎么办？ & 先回粗调；再查棱镜方位（$AB\perp a_1a_2$ 等）、顶点是否朝准直管、台面高矮是否挡光 \\
9 & 怎样判断最小偏向角？ & 慢转台并跟踪像，像\textbf{停住后折返}处即 $\delta_{\min}$；依据是 $\delta(i)$ 在 $i=i'$ 处取极小 \\
10 & 入射光不垂直光栅面会怎样？ & 起作用的缝距变成 $d\cos\gamma$，测出的 $d$ 偏大；$\pm k$ 不再对称、谱线不等高。必须用 $b_1,b_2$ 调垂直 \\
11 & 光栅光谱与棱镜光谱有何不同？哪条线偏折最大？ & 光栅：多级、近似等间距、长波偏折大（红在外）；棱镜：单套、不等间距、短波偏折大（紫在外）。光栅里是 $579.0$ nm 黄线，棱镜里是 $404.7$ nm 紫线 \\
12 & 缝太宽或太窄会怎样？ & 太宽：像成宽带，对准中心带主观误差，两近线并成一条；太窄：太暗，衍射使边缘发虚。取 $0.5$ mm 折中 \\
\bottomrule
\end{tabular}
\end{minipage}
\end{center}

\textbf{两问补充}：\textbf{（讲义 Q7）两条谱线不等高}——竖丝对准斜谱线时，交点位置随高度变，$\theta_k$ 带误差，且 $\pm k$ 不对称，所以要 $b_3$ 调等高。\textbf{（预习 III.9）注意事项}——不碰光栅与棱镜的光学面（指纹擦不掉）、汞灯不频繁开关且不直视、不自己合闸、读数前确认 16 号没动、狭缝与光栅调好后不再碰。

\subhdr{D.4 扩展 3 题（选做）}

\textbf{1. 除反射法外怎么测顶角：自准直法。}用自准直分别让望远镜对准 $AB$ 面与 $AC$ 面的\textbf{法线}（绿十字与上叉丝重合），两次位置之差是\textbf{两法线夹角}；两外法线夹角为 $180^{\circ}-\alpha$，故 $\alpha=180^{\circ}-$ 转角。

\warnbox{\textbf{一个巧合会骗人}：等边棱镜（$\alpha=60^{\circ}$）下，反射法转 $2\alpha=120^{\circ}$，自准直法转 $180^{\circ}-\alpha=120^{\circ}$——\textbf{两种方法转的角一样}。所以只看「转了 $120^{\circ}$」分辨不出用了哪种方法；换 $\alpha=45^{\circ}$ 的棱镜，反射法转 $90^{\circ}$、自准直法转 $135^{\circ}$，差别才出来。作答要写公式，不要只写数字。}

\textbf{2. 除最小偏向角法外怎么测折射率：掠入射法（极限角法）。}让光接近 $90^{\circ}$ 掠射入 $AB$ 面，从 $AC$ 面射出的光有一个\textbf{极限角} $\delta$（再转就消失）。测出它：

\[
n=\sqrt{1+\left(\frac{\sin\delta+\cos\alpha}{\sin\alpha}\right)^{2}}.
\]

推导：掠入射时 $i\to90^{\circ}$，故 $\sin r=1/n$；由 $r+r'=\alpha$ 得 $r'=\alpha-r$；再展开 $n\sin r'=\sin\delta$：

\[
\sin\delta=n\sin(\alpha-r)=n(\sin\alpha\cos r-\cos\alpha\sin r)=\sin\alpha\sqrt{n^{2}-1}-\cos\alpha,
\]

移项平方即得上式。核对：$n=1.5$、$\alpha=60^{\circ}$ 时 $\sin\delta=0.4682$、$\delta=27.92^{\circ}$，代回得 $n=1.500$ $\checkmark$。这个方法不需要找折返点，只要找出射光消失的边界，也是阿贝折射仪的原理基础。

\textbf{3. 分光计还能测什么}：未知波长（已知 $d$ 测 $\varphi_k$，$\lambda=d\sin\varphi_k/k$）；光栅色分辨本领（钠双线能否分开，$R=\lambda/\Delta\lambda$）；棱镜色散曲线（多条谱线的 $\delta_{\min}$ 得 $n(\lambda)$）；棱镜角色散（$D=\Delta\delta/\Delta\lambda$）；多种介质的折射率（最小偏向角法、掠入射法、临界角法）；谱线精细结构（高线数光栅高级次）。

\subhdr{D.5 出门清单}

\begin{itemize}
\item \textbf{预习 3 页的打印件}（p1 题目、p2 思考题、\textbf{p3 数据表}——原文写明「Print it yourself and bring it to class」）。
\item 本书 J1/J2/J3 的表（若不用预习表，就印这一份，别带两张）。
\item 铅笔、橡皮、直尺（作图题要用）、计算器。
\item 校园卡/学生证（签到、借仪器可能要用）。
\end{itemize}

\vspace{6pt}
\begin{center}
{\color{petal}\small\ornstar}\hspace{6pt}{\small\color{inkgray}全书完：J0 调好 $\to$ J1 顶角 $\to$ J2 折射率；J3 衍射角 $\to$ J4 光栅常量与角色散}\hspace{6pt}{\color{petal}\small\ornstar}
\end{center}


\end{document}
